We use the determinant-open quiver of Section~\ref{sec:quiver}
to classify the closed unitary orbits after scalar extension.
We then determine which of these geometric orbits have $F$-rational
points, establish rational orbit uniqueness and transpose stability,
and compute the corresponding normal representations over $F$.
These results provide the input for the distributional arguments in
Section~\ref{sec:unitary-descendants}.

\subsection{The unitary residual action}\label{m2:sec:reduction}

\subsubsection{The group, subgroup, and coordinates}\label{m2:red:pair}

Let $J_{2m}$ be the anti-diagonal matrix with all anti-diagonal
entries equal to one. Write
\[
 G_{2m}(F)=\{g\in\GL_{2m}(E):g^\dagger J_{2m}g
       =\lambda(g)J_{2m},\ \lambda(g)\in F^\times\},
 \qquad g^\dagger=\bar g^{\,t}.
\]
These are the quasi-split unitary similitude groups, with
$G_2=\GU_{1,1}$ and $G_4=\GU_{2,2}$. Put $J=J_2$ and
$\U_2=U_{1,1}=\ker(\lambda:G_2\to\Gm)$ and
$X^*=JX^\dagger J$ for a $2$ by $2$ matrix. The defining equation
of $G_2$ is $h^*h=\lambda(h)I$. We also write
$E^-=\{z\in E:\bar z=-z\}$.
For the Lie algebras, write $\mathfrak{gu}_n=\Lie(G_n)$ and
$\mathfrak u_n=\ker(d\lam:\mathfrak{gu}_n\to F)$.

For $u=(u_{ij})\in G_2$ and $h\in G_2$, define
\[
 \rho(h,u)=\begin{pmatrix}
 u_{11}&0&0&u_{12}\\
 0&h_{11}&h_{12}&0\\
 0&h_{21}&h_{22}&0\\
 u_{21}&0&0&u_{22}
 \end{pmatrix}.
\]
In the notation of the introduction, the unitary pair is
\begin{equation}\label{m2:red:embedding}
 \begin{gathered}
 G_{\mathrm U}=G_4(F)\times G_2(F),\qquad
 H_{\mathrm{eq}}=\{(h,u)\in G_2^2:\lambda(h)=\lambda(u)\},\\
 j(h,u)=(\rho(h,u),h),\qquad H_{\mathrm U}=j(H_{\mathrm{eq}}(F)).
 \end{gathered}
\end{equation}
Order the coordinates as $(1,4,2,3)$, so the outer and middle
planes have form $J\oplus J$ and $\rho(h,u)=u\oplus h$.
This rational permutation preserves ordinary transpose. All subsequent
blocks $x=\left(\begin{smallmatrix}A&B\\C&D\end{smallmatrix}\right)$
use this order; the first plane is the outer plane.

The multiplier map on $G_2(F)$ is surjective: set
\begin{equation}\label{m2:red:amul}
 a_\mu=\diag(1,\mu),\qquad a_\mu^*a_\mu=\mu I
 \quad(\mu\in F^\times).
\end{equation}
The kernel of $H_{\mathrm{eq}}\to\Gm$ is $\U_2^2$.
Thus $H_{\mathrm{eq}}$, like $G_4$ and $G_2$, is connected reductive.
In particular $G_{\mathrm U}$ and $H_{\mathrm U}$ are unimodular,
and $G_{\mathrm U}/H_{\mathrm U}$ has the invariant measure required
in \cref{m2:tools:gk}.

\begin{lemma}\label{m2:red:transpose}
Ordinary transpose preserves $G_2$, $G_4$, and their multipliers.
The anti-involution $\tau(g,z)=(g^t,z^t)$ of $G_{\mathrm U}$ preserves $H_{\mathrm U}$.
\end{lemma}

\begin{proof}
For $J'=J$ or $J\oplus J$, the equation
$g^\dagger J'g=\lambda J'$ implies $gJ'g^\dagger=\lambda J'$.
Conjugating the latter gives $\bar gJ'g^t=\lambda J'$, the defining
equation for $g^t$. Also $\rho(h,u)^t=\rho(h^t,u^t)$, with unchanged
multipliers.
\end{proof}

\subsubsection{The homogeneous base and its fiber}

Let $L=H_{\mathrm{eq}}^2$ act by left and right translation through $j$.
Projection $p(g,z)=z$ has transitive base action
$((h,u),(h',v))z=hzh'^{-1}$, with the continuous section
$z\mapsto((z,a_{\lambda(z)}),(1,1))$ of the orbit map at $1$.
Its fiber is $G_4(F)$; identifying the stabilizer $h=h'$ with
$(u,v,h)$ gives
\begin{equation}\label{m2:red:K}
 K=\{(u,v,h)\in G_2^3:
         \lambda(u)=\lambda(v)=\lambda(h)\}.
\end{equation}
Its action on the fiber is
\begin{equation}\label{m2:red:action}
 (u,v,h)\cdot x=(u\oplus h)x(v\oplus h)^{-1},
 \quad
 (A,B,C,D)\longmapsto
 (uAv^{-1},uBh^{-1},hCv^{-1},hDh^{-1}).
\end{equation}
The multiplier $\lambda(x)$ is unchanged. The kernel of
$K\to\Gm$ is $\U_2^3$, so $K$ is also connected reductive.

\begin{proposition}\label{m2:red:frobenius}
There is a natural isomorphism, compatible with transpose,
\begin{equation}\label{m2:red:distiso}
 \Ddist(G_{\mathrm U})^{H_{\mathrm U}\times H_{\mathrm U}}
   \simeq\Ddist(G_4(F))^{K(F)}.
\end{equation}
\end{proposition}

\begin{proof}
The continuous section constructed above gives the quotient topology on $G_2(F)$.
Since $L$ and $K$ are reductive, their relative modulus is one;
\cref{m2:tools:frobenius} gives \eqref{m2:red:distiso}.
Transpose fixes the base point $1$ and acts on its stabilizer by
\begin{equation}\label{m2:red:alpha}
 \alpha_K(u,v,h)=(v^{-t},u^{-t},h^{-t}),
 \qquad ((u,v,h)\cdot x)^t=\alpha_K(u,v,h)\cdot x^t.
\end{equation}
Its pushforward preserves the positive Haar ray on the base; the
scalar is one because transpose squares to one. Thus Frobenius descent
intertwines total and fiber transpose, as in
\cref{m1:prop:odd-Frobenius}.
\end{proof}

Set
\begin{equation}\label{m2:red:extension}
 \Kt=K\rtimes_{\alpha_K}C_2,
 \qquad \sigma x=x^t,\qquad
 \eps|_K=1,\quad\eps(\sigma)=-1.
\end{equation}
We shall prove
\begin{equation}\label{m2:red:target}
 \Ddist(G_4(F))^{\Kt(F),\eps}=0.
\end{equation}
Taking the odd part of a $K(F)$-invariant distribution shows that
\eqref{m2:red:target} and \cref{m2:red:frobenius} imply the
bi-$H_{\mathrm U}$ distribution hypothesis of \cref{m2:tools:gk}.

\subsubsection{Contragredient duality}

Put $\iota(g,z)=(g^{-t},z^{-t})$, the automorphism associated with $\tau$.

\begin{proposition}\label{m2:red:duality}
For every irreducible admissible smooth complex representation
$\pi$ of $G_{\mathrm U}$, $\pi\circ\iota\simeq\pi^\vee$. Moreover
$\iota(H_{\mathrm U})=H_{\mathrm U}$, and hence
\[
 \dim\Hom_{H_{\mathrm U}}(\pi,\mathbb C)
 =\dim\Hom_{H_{\mathrm U}}(\pi^\vee,\mathbb C).
\]
\end{proposition}

\begin{proof}
Coordinate conjugation $c(v)=\bar v$ is an $F$-linear anti-unitary
involution for our rational symmetric form matrices. Roche--Vinroot's
Theorem~B and its character argument in \S7.2~\cite{RocheVinroot18}
therefore make $\theta_m(g)=\lambda(g)^{-1}\bar g$ a dualizing
involution on $G_{2m}(F)$: $\theta_m(g)$ is conjugate to $g^{-1}$.
The unitary equation gives $g^{-t}=J'\theta_m(g)J'$, where
$J'=J$ or $J\oplus J$ has multiplier one. Thus inverse transpose
is an inner modification of this dualizing involution. The conjugacy
and character argument applies componentwise to the reductive product,
giving $\pi\circ\iota\simeq\pi^\vee$.

Finally $\rho(h,u)^{-t}=\rho(h^{-t},u^{-t})$ and
$\lambda(h^{-t})=\lambda(h)^{-1}$. Hence $\iota$ preserves
$H_{\mathrm U}$, so the invariant linear forms on $\pi\circ\iota$
are exactly those on $\pi$. The contragredient isomorphism gives the
equality of multiplicities.
\end{proof}

\subsubsection{Rational linear algebra of the hyperbolic plane}
\label{m2:red:unitary-pair}

The following rational coordinates will give the transporters and
normal representations for dense orbits.
Choose $\delta\in E^\times$ with $\bar\delta=-\delta$, and set
\[
 d=\delta^2\in F^\times,\qquad P=\diag(1,\delta),\qquad
 J_0=\begin{pmatrix}0&1\\-1&0\end{pmatrix},\qquad
 R=P^{-2}=\diag(1,d^{-1}).
\]
Since $P^\dagger JP=\delta J_0$, the unitary equation becomes an
alternating-form equation. This rational factorization is distinct from
the scalar-extension comparison in \cref{m2:geo:split-bridge}.

\begin{lemma}\label{m2:simconj:factorization}
Every $h\in G_2(F)$ has the form
\begin{equation}\label{m2:simconj:factor}
 h=zPgP^{-1},\qquad z\in E^\times,\quad g\in\GL_2(F),
 \qquad \lambda(h)=N_{E/F}(z)\det g.
\end{equation}
The overlap of two such factorizations is precisely
$(z,g)\sim(z/a,ag)$ for $a\in F^\times$.
Writing $P^{-1}YP=A+\delta B$ with $A,B\in M_2(F)$ defines
an $F$-linear isomorphism
\begin{equation}\label{m2:simconj:Psi}
 \Psi:M_2(E)\longrightarrow M_2(F)^2,
 \quad Y\longmapsto(A,B),\qquad
 \Psi(hYh^{-1})=(gAg^{-1},gBg^{-1}).
\end{equation}
\end{lemma}

\begin{proof}
The identity $g^tJ_0g=(\det g)J_0$ for $g\in\GL_2(F)$ and
$P^\dagger JP=\delta J_0$ show that $PgP^{-1}$ has multiplier
$\det g$; multiplying by $z$ contributes $N(z)$.
Conversely, for $H_0=P^{-1}hP$, comparison of
\[
 \bar H_0^{\,t}J_0H_0=\lambda(h)J_0,
 \qquad H_0^tJ_0H_0=(\det H_0)J_0
\]
gives $\bar H_0=cH_0$, with $c=\lambda(h)/\det H_0$ and
$c\bar c=1$. Hilbert~90 supplies $z$ with $c=\bar z/z$;
then $g=z^{-1}H_0$ is rational. If $zg=z'g'$ with rational
$g,g'$, then $g'g^{-1}=(z/z')I$ forces $z/z'\in F^\times$,
proving the exact overlap.

The $1,\delta$ decomposition is unique, and conjugation by
$h=zPgP^{-1}$ acts on its coefficients by simultaneous $g$-conjugation.
Every $g\in\GL_2(F)$ lifts through $PgP^{-1}$, so $\Psi$
identifies the invariant distribution spaces. The scalar centers
act trivially.
\end{proof}

Since $P^t=P$, the odd coordinate formula is
\begin{align}
 P^{-1}Y^tP
 &=P^{-2}(P^{-1}YP)^tP^2
   =R(A^t+\delta B^t)R^{-1},\notag\\
 \Psi(Y^t)&=(RA^tR^{-1},RB^tR^{-1}).
 \label{m2:simconj:transpose}
\end{align}
Here $R$ belongs to the effective even group $\GL_2(F)$.
The corresponding adjoint identities are as follows.

\begin{lemma}\label{m2:rat:real-pair}
Write $P^{-1}DP=A+\delta B_0$, with
$A,B_0\in\Mat_2(F)$.  Then
\begin{align}
 P^{-1}D^*P&=\operatorname{adj}(A)
                      -\delta\operatorname{adj}(B_0),
               \label{m2:rat:real-adjoint}\\
 P^{-1}D^tP&=R(A^t+\delta B_0^t)R^{-1}.
               \label{m2:rat:real-transpose}
\end{align}
Under this coordinate change
$E\langle D,D^*\rangle=E\langle A,B_0\rangle$.
Every $g\in\GL_2(F)$ gives $PgP^{-1}\in G_2(F)$
with multiplier $\det g$.
\end{lemma}

\begin{proof}
The transported adjoint is $Z\mapsto J_0^{-1}\bar Z^tJ_0$.
The identity
$J_0^{-1}X^tJ_0=\operatorname{adj}(X)=\tr(X)I-X$, linear for
two-by-two matrices, gives \eqref{m2:rat:real-adjoint};
\eqref{m2:simconj:transpose} gives \eqref{m2:rat:real-transpose}.
Adding and subtracting $A+\delta B_0$ and its adjoint, and removing
the scalar trace terms, recovers $A,B_0$ over $E$. This proves the
algebra equality. The multiplier assertion is
\cref{m2:simconj:factorization} with $z=1$.
\end{proof}


\subsection{Geometric orbits and rational existence}\label{m2:geo:section}

We transfer the geometric classification from Section~\ref{sec:quiver},
then determine which parameters admit points in the prescribed Hermitian spaces.

\subsubsection{The block action and the unitary equations}

We use the block coordinates of \cref{m2:red:pair}, with
$\mathcal J=J\oplus J$ and $\mu=\lambda(x)$.
The residual action is \eqref{m2:red:action}; it preserves $\mu$
because its two translation factors have the same multiplier.

\begin{proposition}\label{m2:geo:block-equations}
An invertible block matrix $x$ belongs to $G_4(F)$ with multiplier
$\mu\in F^\times$ if and only if
\begin{align}
 A^\dagger JA+C^\dagger JC&=\mu J,&
 A^\dagger JB+C^\dagger JD&=0,\label{m2:geo:column-equations}\\
 B^\dagger JA+D^\dagger JC&=0,&
 B^\dagger JB+D^\dagger JD&=\mu J.\notag
\end{align}
For such a matrix the equivalent row equations are
\begin{align}
 AJA^\dagger+BJB^\dagger&=\mu J,&
 AJC^\dagger+BJD^\dagger&=0,\label{m2:geo:row-equations}\\
 CJA^\dagger+DJB^\dagger&=0,&
 CJC^\dagger+DJD^\dagger&=\mu J.\notag
\end{align}
At every $F$-point one has $\rk_E B=\rk_E C$.
\end{proposition}

\begin{proof}
The column equations expand $x^\dagger\mathcal Jx=\mu\mathcal J$;
inverting gives $x\mathcal Jx^\dagger=\mu\mathcal J$, hence the row
equations. For the orthogonal block planes $U,W$, restriction of $x$
identifies $\ker B$ with $xW\cap W$ and $\ker C$ with $xU\cap U$.
Since $xW=(xU)^\perp$, one has
$(xW\cap W)^\perp=xU+U$. Comparing dimensions in the four-space
gives $\dim_E(xW\cap W)=\dim_E(xU\cap U)$, proving equal ranks.
\end{proof}

\subsubsection{Passage to the split quiver action}

Fix $\mu\in F^\times$ and write $X_\mu=\lambda^{-1}(\mu)$.
Choose a splitting field $k$ of $E/F$, for example $k=E$ or
$k=\barF$, and order the two factors of
$E\otimes_F k\simeq k\times k$. Conjugation interchanges these
factors; subscripts $1,2$ will denote their matrix components.

\begin{proposition}[Comparison with the split quiver]\label{m2:geo:split-bridge}
There are algebraic isomorphisms over $k$
\[
 (G_4)_k\simeq\GL_4\times\Gm,\qquad
 K_k\simeq\GL_2^3\times\Gm.
\]
Projection to the first matrix component identifies $(X_\mu)_k$
with the determinant-open variety
\[
 \mathscr X=\left\{(A,B,C,D):
 f(A,B,C,D)=\det\begin{pmatrix}A&B\\C&D\end{pmatrix}\ne0\right\}.
\]
The action factors through the three-vertex action
\begin{equation}\label{m2:geo:split-action}
 (p,q,r)\cdot(A,B,C,D)
 =(pAr^{-1},pBq^{-1},qCr^{-1},qDq^{-1}),
 \qquad(p,q,r)=(u_1,h_1,v_1).
\end{equation}
The multiplier torus acts trivially and factors into every geometric
stabilizer. The three-vertex action has its diagonal scalar kernel,
so its effective quotient and determinant character
$\chi(p,q,r)=\det p/\det r$ agree with the linear model. Transpose becomes
\begin{equation}\label{m2:geo:split-transpose}
 (A,B,C,D)\longmapsto(A^t,C^t,B^t,D^t),
\end{equation}
with group automorphism
$(p,q,r,s)\mapsto(r^{-t},q^{-t},p^{-t},s^{-1})$.
\end{proposition}

\begin{proof}
The two-component unitary equation is
$x_2^t\mathcal Jx_1=\mu\mathcal J$. Consequently
\begin{equation}\label{m2:geo:second-component}
 x_2=\mu\mathcal Jx_1^{-t}\mathcal J.
\end{equation}
Thus arbitrary invertible $x_1$ and nonzero $\mu$ determine $x_2$
uniquely and regularly; fixing $\mu$ imposes no condition on
$\det x_1$. Similarly, $u_2=sJu_1^{-t}J$ and the analogous
formulas for $v,h$ identify $K_k$ with $\GL_2^3\times\Gm$.
Substituting their unrestricted first components gives
\eqref{m2:geo:split-action}, independently of $s$, and taking determinants
gives the stated character. The first component of $x^t$ is $x_1^t$;
\eqref{m2:red:alpha} gives the group automorphism, including inversion
of $s$. These are algebraic action and transpose identifications.
\end{proof}

The vertex correspondence is $(U,Q,P)=(\mathcal U,\mathcal W,\mathcal P)$.
The arrows, determinant open, character, and transpose are therefore
those of \cref{m1:eq:X,m1:eq:det-character,m1:eq:residual-action}.
Over $\bar F$, apply the balanced-category and polystability results
\cref{m1:lem:theta-nonnegative,m1:lem:A0-abelian,m1:prop:polystability},
then \cref{m1:lem:stable-factors,m1:prop:rank-exclusion}.
General linear rational normal forms below are taken over $E$;
their Hermitian realization is a separate question.

\subsubsection{Closed families and transpose}

On the dense locus the general linear gauge of \eqref{m1:eq:dense-gauge} is
\begin{equation}\label{m2:geo:dense-gauge}
 \begin{pmatrix}X&I\\I&Y\end{pmatrix},\qquad
 X=B^{-1}AC^{-1},\quad Y=D,
 \qquad(p,q,r)=(B^{-1},I,C).
\end{equation}
The residual action is simultaneous conjugation, with domain
$\det\left(\begin{smallmatrix}X&I\\I&Y\end{smallmatrix}\right)\ne0$.
These coordinates are retained for the Hermitian existence problem.

By \cref{m1:prop:dense-closed}, closedness in the entire determinant
open is geometric semisimplicity of this pair. For rank zero,
\cref{m1:prop:rankzero} gives
\begin{equation}\label{m2:geo:rankzero}
 \begin{pmatrix}I&0\\0&D\end{pmatrix},\qquad
 D\in\GL_2,\qquad\text{closed exactly when $D$ is semisimple.}
\end{equation}
The rational unitary normalization will instead be
$a_\mu\oplus D$, with $D\in G_2(F)$ and $\lambda(D)=\mu$.

\paragraph{Rank-one coordinates.}

For arrow ranks $(1,1)$, \cref{m1:lem:rank11-split,m1:prop:rank11}
give the $E$-rational splitting $\mathrm E\oplus\mathrm H$ of a
closed point, with dimensions $(1,0,1)$ and $(1,2,1)$.
The incidence lines $\im c,\ker b$ are either both nonstable under
$D$, making $\mathrm H$ stable, or both stable and distinct,
giving $\mathrm H=\mathrm R\oplus\mathrm Q$; all other cases are
nonclosed. Hermitian orthogonality will follow from
\cref{m2:rat:selfadjoint-idempotents}.
The Hermitian calculations use the following normal forms.
The decomposable representative \eqref{m1:eq:decomposable-rep}, with
$(x,d,e)=(\zeta,\nu,\omega)$, is
\begin{equation}\label{m2:geo:decomposable-form}
 A=\diag(\zeta,1),\quad B=C=E_{11},\quad
 D=\diag(\nu,\omega),\qquad\omega\ne0,\quad\zeta\nu\ne1.
\end{equation}
Its factors are the identity on the second end coordinates,
$\left(\begin{smallmatrix}\zeta&1\\1&\nu\end{smallmatrix}\right)$
on $(1,1,1)$, and the middle-line loop $\omega$.


For a stable system the cyclic basis $(c,Dc)$ gives
\begin{equation}\label{m2:geo:stable-form}
 A=\diag(a,1),\quad
 B=\begin{pmatrix}b_0&b_1\\0&0\end{pmatrix},\quad C=E_{11},\quad
 D=\begin{pmatrix}0&-d_0\\1&s_0\end{pmatrix},\quad
 b_1=1-ad_0+s_0b_0.
\end{equation}
As in \eqref{m1:eq:b1-normalized}, scaling $\mathcal P_1$ normalizes
$ad_0-s_0b_0+b_1$ to one and scales $a,b_0,b_1$ together.
The target Hermitian form below incorporates this change of basis. Stability is
\begin{equation}\label{m2:geo:stable-pairing}
 H_0=\begin{pmatrix}b_0&b_1\\b_1&s_0b_1-d_0b_0\end{pmatrix}
     =\begin{pmatrix}b\\bD\end{pmatrix},\qquad\det H_0\ne0,
 \quad b=(b_0,b_1).
\end{equation}
These are \eqref{m1:eq:stable-rep-general}--\eqref{m1:eq:H} with
$(d,s)=(d_0,s_0)$, including nilpotent and repeated-root loops;
$d_0$ is independent of $d=\delta^2$ above.

\begin{theorem}[Seven coarse geometric types]\label{m2:geo:classification}
On every fixed similitude fiber the closed geometric $K$-orbits
are precisely the rank-zero semisimple family
\eqref{m2:geo:rankzero}, the dense semisimple-pair family
\eqref{m2:geo:dense-gauge}, and the two rank-one families
\eqref{m2:geo:decomposable-form} and \eqref{m2:geo:stable-form},
subject to the determinant and stability conditions in those formulas.
There are seven coarse geometric types: two rank-zero types
(scalar and distinct-eigenvalue semisimple), three dense types
(absolutely irreducible, two distinct characters, and two equal
characters), and two rank-one types. Each carries continuous
parameters. The nine rational rows arise only after distinguishing
the admissible $F$-forms.
\end{theorem}

\begin{proof}
Over $\bar F$, the comparison in \cref{m2:geo:split-bridge}
identifies the closed orbits with those of the split quiver.
The classification in \cref{m1:thm:closed-orbits} therefore yields
exactly two rank-zero, three dense, and two rank-one types.
\end{proof}

\paragraph{Geometric transpose.}

\begin{lemma}[Symmetric intertwiner]\label{m2:geo:pair-transpose}
Let $(X,Y)$ be a geometrically semisimple pair of two-by-two
matrices over $k_0=F$ or $E$. There is a symmetric invertible
$H\in\GL_2(k_0)$ such that $XH=HX^t$ and $YH=HY^t$.
\end{lemma}

\begin{proof}
Choose $s\in k_0$ with $X+sI$ invertible, then choose $t$ so that
$\det\left(\begin{smallmatrix}X+sI&I\\I&Y+tI\end{smallmatrix}\right)\ne0$:
as a polynomial in $t$ it has degree two and leading coefficient
$\det(X+sI)\ne0$. Scalar translations preserve the generated algebra
and semisimplicity. Thus \cref{m1:lem:dense-transpose} applies to this
closed dense point. Subtracting the scalar terms from its intertwiner
equations and taking $H=H_{\mathrm{lin}}^{-1}$ gives the required
symmetric invertible intertwiner in our convention.
\end{proof}

\begin{proposition}[Geometric transpose stability]
\label{m2:geo:transpose}
Every individual closed geometric orbit in
\cref{m2:geo:classification} is preserved by ordinary transpose.
A split-quiver point over $E$ with closed geometric orbit has
an $E$-linear general linear transporter from its transpose to itself.
\end{proposition}

\begin{proof}
The row formulas of \cref{m1:cor:transpose-stable}, transported
through \cref{m2:geo:split-bridge}, apply over $E$ to the general
linear transporter and over $\bar F$ to geometric orbit stability. The dense
transporter is $(H,H,H)$; in rank zero apply the symmetric
intertwiner to $(D,I)$ and take $(I,H,I)$. The decomposable
representative is symmetric, while the stable system has
\begin{equation}\label{m2:geo:stable-transporter}
 (I,H_0^{-1},I)\cdot(A^t,C^t,B^t,D^t)=(A,B,C,D),
\end{equation}
in the $(p,q,r)$ order, with $H_0$ as in
\eqref{m2:geo:stable-pairing}.
The general linear normalizations transport this construction to
an $E$-linear transporter at every point. Rational unitary transpose
stability, including common-multiplier compatibility, is proved in
\cref{m2:rat:transpose}.
\end{proof}

\subsubsection{Hermitian planes and rank-zero existence}

Write
\[
 Q_F=F^\times/\N_{E/F}(E^\times).
\]
This group has order two.  Every element has square one, because
$s^2=\N_{E/F}(s)$ for $s\in F^\times$.  For a nondegenerate
two-dimensional Hermitian matrix $H$ define its relative
discriminant, relative to the hyperbolic matrix $J$, by
\begin{equation}\label{m2:rat:discriminant}
 \operatorname{disc}_J(H)=[-\det H]\in Q_F.
\end{equation}
The minus sign divides by $\det J=-1$.

\begin{lemma}\label{m2:rat:hermitian-planes}
A nondegenerate Hermitian plane over $E/F$ is hyperbolic if and
only if $-\det H\in\N_{E/F}(E^\times)$.  Two Hermitian planes
over the local field $F$ are isometric if and only if their
determinants have the same class in $Q_F$.  Scaling a plane by
an element of $F^\times$ does not change this class.
\end{lemma}

\begin{proof}
Hermitian orthogonalization diagonalizes the form as
$\diag(a,b)$, with $a,b\in F^\times$.  It has a nonzero
isotropic vector exactly when $-b/a$ is a norm.  This is
equivalent to $-ab$ being a norm, since $a^2$ is a norm.
An isotropic nondegenerate plane is hyperbolic: choose an
isotropic vector $e$, choose $f$ with $\langle e,f\rangle=1$,
and replace $f$ by
$f-\frac12\langle f,f\rangle e$ to obtain an isotropic
partner. This proves the first assertion over any quadratic
extension of characteristic different from two.

For the second assertion, put $d=-b/a$ and form the quaternion
algebra $\mathscr B=E\oplus Ej$, with $j^2=d$ and
$jz=\bar z j$ for $z\in E$. Its reduced norm satisfies
\[
 a\operatorname{Nrd}(x+yj)
 =a\N_{E/F}(x)+b\N_{E/F}(y).
\]
Over a non-Archimedean local field the reduced norm
$\operatorname{Nrd}:\mathscr B^\times\to F^\times$ is
surjective~\cite[Lemma~13.4.9]{Voight21}. Thus the plane represents $1$.
Orthogonalizing at such a vector writes it as $\diag(1,c)$,
where $[c]=[\det H]$ in $Q_F$. Two such forms are isometric
when their second coefficients differ by a norm, and the
converse follows from the change-of-basis formula for
determinants. This proves the plane case of
\cite[Theorem~3.1]{Jacobowitz62}. Finally
$\det(sH)=s^2\det H$ and $s^2=\N_{E/F}(s)$.
\end{proof}

We express the hyperbolicity conditions using \eqref{m2:rat:discriminant}.

Fix $\mu\in F^\times$ and put $a_\mu=\diag(1,\mu)$.
If $B=C=0$, the unitary equations say $A,D\in G_2(F)$ with
multiplier $\mu$.  The element $(a_\mu A^{-1},1,1)\in K(F)$
normalizes $A$ to $a_\mu$.  Hence all rank-zero representatives
are
\begin{equation}\label{m2:rat:rankzero-representative}
 x_{\mu,D}=a_\mu\oplus D,
 \qquad D\in G_2(F),\quad\lam(D)=\mu.
\end{equation}
Two such representatives lie in the same $K(F)$-orbit exactly
when their $D$-blocks are conjugate in $G_2(F)$: one
implication follows from the last block of
\eqref{m2:red:action}. Conversely, if $hDh^{-1}=D'$, take
$v=a_{\lambda(h)}$ and $u=a_\mu va_\mu^{-1}$. Both have
multiplier $\lambda(h)$, and $(u,v,h)$ realizes the required
conjugacy in $K(F)$. The stabilizer embedding is
\begin{equation}\label{m2:rat:rankzero-stabilizer}
 K_{x_{\mu,D}}
 =\{(a_\mu v a_\mu^{-1},v,h):
       hD=Dh,\ \lam(v)=\lam(h)\}
 \simeq G_2\times_{\lam}Z_{G_2}(D).
\end{equation}
Closedness is equivalent to semisimplicity of $D$ by
\cref{m2:geo:classification}.  If $D=zI$, admissibility is
$\N_{E/F}(z)=\mu$ and the stabilizer is
$G_2\times_\lam G_2$.  If $D$ is regular semisimple,
$M=E[D]$ is quadratic etale over $E$.  Since
$D^*=\mu D^{-1}$, it is stable under $*$, and
\[
 Z_{G_2}(D)=T_{M,*}.
\]
The three algebra forms in \cref{m2:rat:torus-cohomology}
distinguish the split and the two anisotropic possibilities
for its quotient by the center.

\subsubsection{The dense Hermitian domain}

For a matrix $D\in\Mat_2(E)$ and a fixed $\mu\in F^\times$
define
\begin{align}
 M_D&=\mu I-D^*D,& Q_D&=JM_D,\label{m2:rat:dense-Gram}\\
 X_D&=-M_D^{-1}D^*,&
 R_D&=J+X_D^\dagger Q_DX_D,\notag
\end{align}
whenever $M_D$ is invertible.  The matrices $Q_D$ and $R_D$
are Hermitian.  Let $\Omega_\mu$ consist of those $D$ for
which $M_D$ and $R_D$ are invertible and
\begin{equation}\label{m2:rat:dense-domain}
 [-\det Q_D]=[-\det R_D]=1\quad\hbox{in }Q_F.
\end{equation}
These conditions say that $Q_D$ is isometric to $J$
and $R_D$ is isometric to $\mu J$; the latter is again
hyperbolic.  Thus \eqref{m2:rat:dense-domain} includes all
Hermitian-form obstructions to the choices below.

\begin{proposition}\label{m2:rat:dense-param}
For $D\in\Omega_\mu$, choose invertible $B,C$ with
\begin{equation}\label{m2:rat:dense-factorizations}
 B^\dagger JB=Q_D,\qquad C^\dagger R_DC=\mu J.
\end{equation}
Then
\begin{equation}\label{m2:rat:dense-representative}
 g_{\mu,D;B,C}=
 \begin{pmatrix}BX_DC&B\\C&D\end{pmatrix}
\end{equation}
belongs to $G_4(F)$ with multiplier $\mu$ and has rank
$(2,2)$.  Every rank-$(2,2)$ point is obtained in this way.
Its orbit depends precisely on the $G_2(F)$-conjugacy class
of $D$ in $\Omega_\mu$.  Its stabilizer is
\begin{equation}\label{m2:rat:dense-stabilizer}
 K_g=\{(BhB^{-1},C^{-1}hC,h):h\in Z_{G_2}(D)\}
 \simeq Z_{G_2}(D).
\end{equation}
\end{proposition}

\begin{proof}
The definitions give
\[
 Q_DX_D=-JD^*,\qquad X_D^\dagger Q_D=-JD.
\]
Hence the off-diagonal equation is
\[
 A^\dagger JB+C^\dagger JD
 =C^\dagger(X_D^\dagger Q_D+JD)=0.
\]
The last diagonal equation is
$B^\dagger JB+D^\dagger JD=\mu J$ by the definition of
$Q_D$, and the first is
\[
 A^\dagger JA+C^\dagger JC
 =C^\dagger(X_D^\dagger Q_DX_D+J)C=\mu J.
\]
These verify all equations of \cref{m2:geo:block-equations}.
Since the right-hand side is nondegenerate, the resulting
matrix is invertible.

Conversely, let the point have invertible $B,C$ and set
$X=B^{-1}AC^{-1}$.  The last diagonal equation gives
$B^\dagger JB=Q_D$, so $M_D$ is invertible.  The
off-diagonal equation gives $X^\dagger Q_D=-JD$, whence
$X=X_D$.  The first diagonal equation then gives the second
factorization in \eqref{m2:rat:dense-factorizations}.  Thus
$D\in\Omega_\mu$ and the point has the claimed form.

For the same $D$, let $(B',C')$ be another pair of choices.
Then $u=B'B^{-1}$ and $v=C'^{-1}C$ are unitary: the first
identity follows from their common Gram matrix $Q_D$, and
the second from their common $R_D$-Gram matrix $\mu J$.
The element $(u,v,1)$ carries the original representative
to the new one.  Under an arbitrary residual element, $D$
changes by $h$-conjugation, and the construction is equivariant.
This proves the orbit assertion.

The block equations for a stabilizing triple force
$u=BhB^{-1}$, $v=C^{-1}hC$, and $hD=Dh$.
Conversely, if $h\in G_2$ commutes with $D$, taking adjoints
shows that it commutes with $D^*$ as well.  It therefore
commutes with $X_D$ and preserves the two Gram forms up to
its multiplier.  For example,
$h^\dagger Q_Dh=\lam(h)Q_D$.  This shows that
$BhB^{-1}$ belongs to $G_2$ with multiplier $\lam(h)$.
Similarly $h^\dagger R_Dh=\lam(h)R_D$, and the second
factorization proves the same assertion for $C^{-1}hC$.
The remaining $A$-equation follows from $hX_D=X_Dh$.
\end{proof}

The general linear gauge of \eqref{m2:geo:dense-gauge} at this
point is $(X_D,D)$.  Besides $X_D\in E[D,D^*]$, one has
\begin{equation}\label{m2:rat:dense-algebra}
 D^*=-\mu X_D(I-DX_D)^{-1},
 \qquad E\langle X_D,D\rangle=E\langle D,D^*\rangle.
\end{equation}
For the first identity use
$X_D=-D^*(\mu I-DD^*)^{-1}$ and
$I-DX_D=\mu(\mu I-DD^*)^{-1}$.
All inverses in this calculation exist because
$\mu I-D^*D$ is invertible.  Thus
\cref{m1:prop:dense-closed} says that the dense orbit is closed
exactly when the natural module for $E\langle D,D^*\rangle$
is semisimple.  Characteristic zero ensures that
semisimplicity is preserved by algebraic extension.

The resulting stabilizers are as follows.  If the algebra is
$\Mat_2(E)$, its commutant is $E$, so the stabilizer is
$\Res_{E/F}\Gm$.  If $D$ is scalar, the stabilizer is
$G_2$.  Otherwise the semisimple commutative algebra is
$M=E[D]$, a quadratic etale $E$-algebra stable under $*$,
and the stabilizer is $T_{M,*}$. The algebra forms in
\cref{m2:rat:torus-cohomology} give its split or anisotropic quotient
by the center.

\subsubsection{Rank-one orthogonality and existence}



The two rank-one normal forms of \cref{m2:geo:classification} are
general linear normal forms over $E$.  We next describe
exactly which Hermitian structures on them give rational
points in the prescribed hyperbolic spaces.  The description
also proves the stabilizer assertions uniformly in their
parameters.

For an $E$-quiver representation with block isomorphism $g$,
let $\mathscr A_g$ be its endomorphism algebra.  In the
$(u,v,h)$ order it consists of triples satisfying
\begin{equation}\label{m2:rat:endo-equations}
 uA=Av,\qquad uB=Bh,\qquad hC=Cv,\qquad hD=Dh.
\end{equation}
Invertibility is not imposed on these endomorphisms.
For a decomposable closed object, the stable-factor
decomposition is
\[
 \mathrm E\oplus\mathrm R\oplus\mathrm Q,
 \qquad
 (1,0,1),\ (1,1,1),\ (0,1,0),
\]
and $\mathscr A_g=E^3$.  For a stable system it is
$\mathrm E\oplus\mathrm H$, of dimensions
$(1,0,1),(1,2,1)$, and $\mathscr A_g=E^2$.
These scalar endomorphism algebras and their vertex embeddings
are the stable-factor data of \cref{m1:prop:rank11,m1:eq:decomposable-stabilizer,m1:eq:stable-stabilizer}, applied over $E$.
These are the factor and endomorphism structures over $E$. The next
step incorporates the Hermitian adjoint and orthogonality to determine
the unitary stabilizers over $F$.

\begin{lemma}\label{m2:rat:selfadjoint-idempotents}
Suppose such a rank-one object is equipped with nondegenerate
Hermitian forms at its vertices and its block map is a
similitude of multiplier $\mu$.  Its endomorphism algebra
is stable under the componentwise Hermitian adjoint, and
every primitive stable-factor idempotent is selfadjoint.
Consequently the stable-factor decompositions at each
vertex are orthogonal decompositions into nondegenerate
spaces.
\end{lemma}

\begin{proof}
Write an endomorphism as $Lg=gR$, with
$L=\diag(u,h)$ and $R=\diag(v,h)$.  The adjoints use the
appropriate target and source vertex forms.  Since
$g^*=\mu g^{-1}$, taking adjoints and multiplying by $g$
gives
\[
 R^*g^{-1}=g^{-1}L^*,\qquad L^*g=gR^*.
\]
Thus $(u^*,v^*,h^*)$ is again an endomorphism.  The induced
semilinear involution of $E^3$ or $E^2$ permutes primitive
idempotents.  Adjoint preserves the rank at every vertex,
whereas the dimension vectors of these idempotents are
distinct.  It can therefore permute none of them.

For a selfadjoint idempotent $e$, its image is orthogonal
to its kernel. Their direct sum is the entire
nondegenerate vertex space, so both restrictions are
nondegenerate. Every primitive idempotent has this property.
\end{proof}

\begin{proposition}\label{m2:rat:rankone-stabilizers}
Put
\[
 T_m=\{(z_0,\ldots,z_{m-1})\in(E^\times)^m:
          \N(z_0)=\cdots=\N(z_{m-1})\}.
\]
The rational stabilizer of every decomposable rank-one
point is $T_3$, and that of every stable-system point is
$T_2$.  In the stable-factor coordinates their vertex
representations are
\begin{equation}\label{m2:rat:rankone-embeddings}
 \begin{array}{c|ccc}
 &u\text{ on }\mathcal P&v\text{ on }\mathcal U
                       &h\text{ on }\mathcal W\\ \hline
 T_3:(r,a_1,b_2)&\diag(r,a_1)&\diag(r,a_1)&\diag(r,b_2)\\
 T_2:(r,a_1)&\diag(r,a_1)&\diag(r,a_1)&rI_2.
 \end{array}
\end{equation}
The multiplier is the common norm.  All these stabilizers
are connected tori, of dimensions four and three respectively.
\end{proposition}

\begin{proof}
The automorphisms are exactly the invertible scalar
endomorphisms of the stable factors.  By
\cref{m2:rat:selfadjoint-idempotents} each factor is
nondegenerate at every vertex where it occurs.  Scalar
multiplication on it has multiplier the norm of that
scalar.  The requirement that all three vertex
similitudes have a common multiplier therefore says
exactly that all the scalar norms agree, with the
embeddings in \eqref{m2:rat:rankone-embeddings}.
Over $\barF$ the equations are equality of the products
of two independent coordinates for each scalar.  Choosing
their common product and one coordinate in each pair
identifies $T_m$ with a split torus of dimension $m+1$
after base change.  In particular there are no additional
finite components.
\end{proof}

\paragraph{Hermitian parameter domains.}

We describe the rank-one parameter domains by equations and
norm conditions.  Let $g_0$ be one of the
general linear normal forms of
\eqref{m2:geo:decomposable-form} or \eqref{m2:geo:stable-form},
with entries in $E$.  Choose $\mu\in F^\times$ and
Hermitian matrices $\mathsf H_U,\mathsf H_P,\mathsf H_W$ satisfying
\begin{equation}\label{m2:rat:rankone-Hermitian}
 g_0^\dagger\diag(\mathsf H_P,\mathsf H_W)g_0
       =\mu\diag(\mathsf H_U,\mathsf H_W),
 \qquad
 \det \mathsf H_i\ne0,\quad[-\det \mathsf H_i]=1\text{ in }Q_F.
\end{equation}
The same form $\mathsf H_W$ occurs on the source and target
middle vertex.  The final three conditions specify that
all vertex planes are the prescribed hyperbolic planes.

For the decomposable domain $\mathscr R_{\mathrm{dec}}$,
take $g_0$ as in \eqref{m2:geo:decomposable-form} and write
\begin{equation}\label{m2:rat:decomposable-forms}
 \mathsf H_U=\diag(u_r,u_a),\qquad
 \mathsf H_P=\diag(p_r,p_a),\qquad
 \mathsf H_W=\diag(w_r,w_b),
\end{equation}
where all six entries belong to $F^\times$.  Orthogonality
in \cref{m2:rat:selfadjoint-idempotents} shows that these
diagonal forms lose no solutions.  The block matrix
equation in \eqref{m2:rat:rankone-Hermitian} is exactly
\begin{align}
 p_a&=\mu u_a,& \N(\omega)&=\mu,
       \label{m2:rat:decomposable-domain}\\
 p_r\N(\zeta)+w_r&=\mu u_r,&
 p_r\bar\zeta+w_r\nu&=0,\notag\\
 p_r+w_r\N(\nu)&=\mu w_r.&&\notag
\end{align}
To verify this, restrict to the $(1,1,1)$ stable factor.
Its matrix $\left(\begin{smallmatrix}\zeta&1\\1&\nu\end{smallmatrix}\right)$
is a similitude from $\diag(u_r,w_r)$ to
$\diag(p_r,w_r)$, giving the last three equations and
their conjugates.  The other two factors give the first
two equations.  The hyperbolicity requirements are
\begin{equation}\label{m2:rat:decomposable-norms}
 [-u_ru_a]=[-p_rp_a]=[-w_rw_b]=1\quad\text{in }Q_F.
\end{equation}
Together with $\omega\ne0$ and $\zeta\nu\ne1$,
\eqref{m2:rat:decomposable-domain} and
\eqref{m2:rat:decomposable-norms} define the whole domain.
In particular a decomposable rational point can occur
only when $\mu$ is a norm.  This restriction follows
from the scalar factor on the same middle line in both
source and target.

For the stable-system domain $\mathscr R_{\mathrm{st}}$,
take \eqref{m2:geo:stable-form}, require $\det H_0\ne0$,
and write
\begin{equation}\label{m2:rat:stable-forms}
 \mathsf H_U=\diag(u_r,u_a),\qquad \mathsf H_P=\diag(p_r,p_a),
 \qquad
 \mathsf H_W=\begin{pmatrix}w_{11}&w_{12}\\
                    \bar w_{12}&w_{22}\end{pmatrix}.
\end{equation}
Here $u_r,u_a,p_r,p_a\in F^\times$,
$w_{11},w_{22}\in F$, and $w_{12}\in E$.  The middle
factor is the stable two-dimensional space.  With
$b=(b_0,b_1)$, equation \eqref{m2:rat:rankone-Hermitian}
becomes
\begin{align}
 p_a&=\mu u_a,\label{m2:rat:stable-domain}\\
 p_r\N(a)+w_{11}&=\mu u_r,\notag\\
 \bar a\,p_rb+e_1^\dagger \mathsf H_WD&=0,\notag\\
 p_rb^\dagger b+D^\dagger \mathsf H_WD&=\mu \mathsf H_W.\notag
\end{align}
The remaining off-diagonal block is the adjoint of the
third equation.  These equations, the nonzero determinant
and norm conditions of \eqref{m2:rat:rankone-Hermitian},
and the stability inequality $\det H_0\ne0$ are the
complete restrictions.  The determinant normalization
of the $(1,2,1)$ block also gives the useful consequence
$p_r=\mu^3u_r$.

\begin{proposition}\label{m2:rat:rankone-param}
For either of these domains choose $E$-linear isometries
$b_U,b_P,b_W$ such that $\mathsf H_i=b_i^\dagger Jb_i$ and put
\begin{equation}\label{m2:rat:rankone-transport-forms}
 g=(b_P\oplus b_W)g_0(b_U\oplus b_W)^{-1}.
\end{equation}
This gives a rank-one closed point of $G_4(F)$ with
multiplier $\mu$.  Conversely, every rank-one closed
$F$-point is obtained from one of these domains.
\end{proposition}

\begin{proof}
The isometries exist by the hyperbolicity conditions.
Substitution of $\mathsf H_i=b_i^\dagger Jb_i$ into
\eqref{m2:rat:rankone-Hermitian} proves
$g^\dagger\mathcal Jg=\mu\mathcal J$ for
\eqref{m2:rat:rankone-transport-forms}.  Changes of basis
preserve the arrow ranks and the geometric closedness
of the normal form.

Conversely, regard a rational unitary point as a quiver
representation over $E$.  The classification and the
linear retraction in \cref{m1:lem:rank11-split} give
an $E$-change of basis to the appropriate general linear
normal form. The original vertex forms $J$, transported along
these changes of basis, remain hyperbolic and satisfy
\eqref{m2:rat:rankone-Hermitian}.
\Cref{m2:rat:selfadjoint-idempotents} makes their stable-factor
decompositions orthogonal, yielding
\eqref{m2:rat:decomposable-forms} or \eqref{m2:rat:stable-forms}.
Expanding the block equation gives exactly the listed
domain equations.  In the stable case, the determinant-one normalization at
$\mathcal P_1$ is a change of basis accounted for by the transported
$\mathsf H_P$.
\end{proof}

These conditions describe the existence domains;
Section~\ref{m2:rat:section} establishes uniqueness of rational orbits
within each admitted geometric orbit.

\subsection{Rational orbit uniqueness and transpose}\label{m2:rat:section}

The existence domains in Section~\ref{m2:geo:section} supply the rational base points. We now compute the stabilizer-cohomology kernels at every admitted point.

\begin{theorem}[Rational closed-orbit uniqueness]\label{m2:rat:classification}
For every closed geometric $K$-orbit $\mathscr O$ containing an
$F$-point,
\begin{equation}\label{m2:rat:unique-orbit}
 K(F)\backslash\mathscr O(F)=\{*\}.
\end{equation}
The analytically closed rational orbits form nine families:
three rank-zero families (scalar, regular split-adjoint, and regular
anisotropic-adjoint), four dense families (absolutely irreducible,
scalar, split-adjoint toral, and anisotropic-adjoint toral), and the
decomposable and stable-system rank-one families. Their stabilizers
are given by \eqref{m2:rat:rankzero-stabilizer},
\eqref{m2:rat:dense-stabilizer}, and
\cref{m2:rat:rankone-stabilizers}; the full list is recorded with the
normal representations in \cref{m2:nor:classification}.
The parameter domains are
the rank-zero condition \eqref{m2:rat:rankzero-representative} with
$D$ semisimple, the dense domain \eqref{m2:rat:dense-domain} with
$E\langle D,D^*\rangle$ acting semisimply, and the complete
decomposable and stable-system domains of
\cref{m2:rat:rankone-param}. All their determinant, multiplier, and
hyperbolicity conditions are retained.
\end{theorem}


\subsubsection{Orbit cocycles and ambient cohomology}

\begin{proposition}[The pointed-set orbit correspondence]
\label{m2:rat:orbit-correspondence}
Fix a closed geometric orbit $\mathscr O$ containing
$x\in G_4(F)$, and let $K_x$ be its algebraic stabilizer over $F$.
Then there is a natural pointed-set bijection
\begin{equation}\label{m2:rat:orbit-kernel}
 K(F)\backslash\mathscr O(F)
 \simeq\ker\bigl[H^1(F,K_x)\longrightarrow H^1(F,K)\bigr],
\end{equation}
with the orbit of $x$ corresponding to the neutral class.
\end{proposition}

\begin{proof}
Put $\Gamma=\Gal(\barF/F)$. For $y=kx\in\mathscr O(F)$,
the elements $c_\gamma=k^{-1}\gamma(k)$ lie in $K_x(\barF)$
and satisfy $c_{\gamma\eta}=c_\gamma\gamma(c_\eta)$.
Their continuous cocycle class has neutral image in $H^1(F,K)$.
Changing the lift to $kt$, with $t\in K_x(\barF)$, replaces it by
$t^{-1}c_\gamma\gamma(t)$; left multiplication by $K(F)$ leaves
it unchanged. This defines the map on rational orbits.

Conversely, a stabilizer cocycle with neutral ambient image is
$k^{-1}\gamma(k)$, and $\gamma(kx)=kc_\gamma x=kx$ gives a
rational point. If lifts $k,l$ give cohomologous stabilizer cocycles,
choose $t\in K_x(\barF)$ with
$l^{-1}\gamma(l)=(kt)^{-1}\gamma(kt)$. Then $l(kt)^{-1}\in K(F)$
sends $kx$ to $lx$, proving injectivity. Finally $k=1$ corresponds
to the neutral class.
\end{proof}

The domains of Section~\ref{m2:geo:section} decide nonemptiness.
Every kernel calculation below starts with an admitted rational point.

We use nonabelian Galois cohomology with its distinguished neutral
class, as in \cite{SerreGC}.  Descent of the defining Hermitian
space identifies $H^1(F,U_{1,1})$ with the isometry classes of
nondegenerate Hermitian planes: after forgetting the form, a
cocycle is trivial by matrix Hilbert~90, and its descended form
is well defined up to an $E$-linear isometry. Conversely every
such form becomes isometric to $J$ over $\bar F$ and hence
defines a unitary cocycle. By \cref{m2:rat:hermitian-planes},
relative discriminant therefore gives
\[
 H^1(F,U_{1,1})\simeq Q_F.
\]
The multiplier sequence
\begin{equation}\label{m2:rat:ambient-sequence}
 1\longrightarrow U_{1,1}^3\longrightarrow K
 \xrightarrow{\lam}\Gm\longrightarrow1
\end{equation}
has the rational section
$s\mapsto(a_s,a_s,a_s)$, $a_s=\diag(1,s)$.
To see the resulting quotient directly, the multiplier of a
$K$-cocycle is a $\Gm$-cocycle. Hilbert~90 and the section
give a cohomologous cocycle in $U_{1,1}^3$. If two such cocycles
are cohomologous by $b\in K(\bar F)$, then $\lambda(b)$ is
Galois-fixed. Writing $b$ as a unitary triple times the section
of $s=\lambda(b)\in F^\times$ shows that their classes in
$H^1(F,U_{1,1}^3)$ differ precisely by the $F^\times$-action
induced by this section. Thus $H^1(F,K)$ is the quotient of
$Q_F^3$ by that action. On a Hermitian form it is scalar
similarity, up to change of basis.  Its determinant contribution in
each two-dimensional factor is $s^2$, hence a norm.  The action
on all three discriminants is therefore trivial, and
\begin{equation}\label{m2:rat:H1K}
 H^1(F,K)\simeq Q_F^3.
\end{equation}
The same argument applied to
\[
 1\longrightarrow U_{1,1}\longrightarrow G_2
 \xrightarrow{\lambda}\Gm\longrightarrow1,
 \qquad s\longmapsto a_s,
\]
gives $H^1(F,G_2)\simeq Q_F$.

\subsubsection{Regular involutive tori}



The same tori occur in the regular rank-zero and dense toral
families.  Let $M$ be a quadratic etale $E$-algebra with a
semilinear involution $*$ extending conjugation on $E$, and put
\[
 L=M^{*=1},\qquad
 T_{M,*}=\{r\in M^\times:r^*r\in F^\times\}.
\]
Then $L$ is a quadratic etale $F$-algebra and $M=L\otimes_FE$.
The multiplier of $r$ is $r^*r$.  These definitions refer to
algebraic groups over $F$, as well as to their rational points.

There are three possible involutive-algebra forms:
\begin{enumerate}
 \item $M=E\times E$ and
 $(r_1,r_2)^*=(\bar r_2,\bar r_1)$; then $L\simeq E$;
 \item $M=E\times E$ and
 $(r_1,r_2)^*=(\bar r_1,\bar r_2)$; then $L=F\times F$;
 \item $M$ is a field; then $L/F$ is quadratic and distinct
 from $E$, and $M=EL$ is biquadratic over $F$.
\end{enumerate}
To see the last assertion, the fixed field of $*$ has degree
two over $F$, is not $E$, and together with $E$ generates $M$.
The split cases follow by asking whether $*$ preserves or
interchanges the primitive idempotents.

\begin{proposition}\label{m2:rat:torus-cohomology}
There is a natural identification
\begin{equation}\label{m2:rat:H1T}
 H^1(F,T_{M,*})
 \simeq L^\times/\bigl(F^\times\N_{M/L}(M^\times)\bigr).
\end{equation}
Suppose $T_{M,*}$ is embedded in $G_2$ through a rank-one
Hermitian $M$-module whose underlying $E/F$-Hermitian plane
is hyperbolic. The map to $H^1(F,G_2)$ forgets the rank-one
$M$-module structure and retains the underlying Hermitian plane
up to scalar similarity. By \cref{m2:rat:hermitian-planes}, scalar
similarity does not change its isometry class in this dimension.
The induced map is the relative determinant map
\begin{equation}\label{m2:rat:Delta}
 \Delta:[c]\longmapsto[\N_{L/F}(c)]\in Q_F.
\end{equation}
For each of the three involutive-algebra forms, $\Delta$ is
injective.
\end{proposition}

\begin{proof}
The norm morphism gives an exact sequence of tori
\[
 1\longrightarrow T_{M,*}
 \longrightarrow\Res_{M/F}\Gm
 \xrightarrow{\ r\mapsto r^*r\bmod F^\times\ }
 (\Res_{L/F}\Gm)/\Gm\longrightarrow1.
\]
Exactness follows after base change from the corresponding
sequence of products of multiplicative groups. Hilbert~90
annihilates the first cohomology of the middle torus.
Moreover
\[
 ((\Res_{L/F}\Gm)/\Gm)(F)=L^\times/F^\times,
\]
because $H^1(F,\Gm)=1$.  The connecting map consequently gives
\eqref{m2:rat:H1T}.

To calculate the map induced by the two-dimensional
representation, identify its rank-one $M$-module with $M$
and write the given Hermitian form as
\[
 \langle x,y\rangle_{c_0}
   =\Tr_{M/E}(c_0xy^*),\qquad c_0\in L^\times.
\]
Indeed $\ell(z)=\langle z,1\rangle$ is $E$-linear, and
the adjoint identity gives
$\langle x,y\rangle=\ell(xy^*)$. Nondegeneracy of the
trace pairing writes $\ell(z)=\Tr_{M/E}(c_0z)$ uniquely.
Hermitian symmetry then gives $c_0^*=c_0$, and
nondegeneracy of the form gives $c_0\in L^\times$.
A cohomology class represented by
$c$ changes this coefficient to $c_0c$, modulo a common
$F$-scalar and an $M/L$ norm.  The determinant ratio is
$\N_{M/E}(c)=\N_{L/F}(c)$.  This proves
\eqref{m2:rat:Delta}; in particular the map is independent of
the initial coefficient $c_0$.

We evaluate all three involutive-algebra cases explicitly.
For the componentwise involution on $M=E\times E$, one has
$L=F\times F$ and
\[
 H^1(F,T_{M,*})=Q_F^2/\diag(Q_F),\qquad
 \Delta([c_1,c_2])=[c_1c_2].
\]
Since $Q_F$ has exponent two, the product map has exactly the
diagonal kernel and induces an isomorphism to $Q_F$.

If $L/F$ is a field distinct from $E$, write
$E=F(\sqrt a)$.  Compatibility of quadratic norm characters
with the field norm, or equivalently the projection formula
for Hilbert symbols, is
\begin{equation}\label{m2:rat:norm-projection}
 (c,a)_L=(\N_{L/F}(c),a)_F.
\end{equation}
Here $(c,a)_L=1$ means exactly that $c$ is a norm from
$L(\sqrt a)=M$.  Formula \eqref{m2:rat:norm-projection} follows
from the norm functoriality of local reciprocity
\cite[Proposition~4.3.11]{KedlayaCFT}: under
$\Gal(EL/L)\simeq\Gal(E/F)$, the restriction to $E$
of $\phi_F(\N_{L/F}c)$ equals the restriction to
$EL$ of $\phi_L(c)$.  Evaluation at the nontrivial
quadratic character yields \eqref{m2:rat:norm-projection}.
Here $\phi_F$ and $\phi_L$ denote the respective
local reciprocity maps.
In particular it gives
\[
 c\in\N_{M/L}(M^\times)
 \quad\Longleftrightarrow\quad
 \N_{L/F}(c)\in\N_{E/F}(E^\times).
\]
For $c\in F^\times$, its norm from $L$ is $c^2$, so every
$F$-scalar already lies in $\N_{M/L}(M^\times)$.  Thus the
kernel of \eqref{m2:rat:Delta} is trivial.  The map is also
onto $Q_F$: the norm subgroups for the two distinct quadratic
extensions $E/F$ and $L/F$ are distinct subgroups of index
two by local reciprocity
\cite[Theorem~4.1.2]{KedlayaCFT}, so
$\N_{L/F}(L^\times)$ is not contained in
$\N_{E/F}(E^\times)$.

Finally, for the exchanged involution on $M=E\times E$,
one has $L\simeq E$. The algebra $M/L$ is split, its norm
is surjective, and the group in \eqref{m2:rat:H1T} is trivial.
This completes all cases.
\end{proof}

\paragraph{The effective root torus.}

For the distributional arguments we also need the quotient of
$T_{M,*}$ by its scalar subgroup.

It is useful to distinguish $L$ from the quadratic algebra that
splits the effective root torus.  Let $\iota$ be the nontrivial
$E$-automorphism of $M$ and set
$K_T=M^{\iota *=1}$.

\begin{lemma}\label{m2:rat:adjoint-torus}
There is an exact sequence of algebraic $F$-tori
\begin{equation}\label{m2:rat:root-sequence}
 1\longrightarrow\Res_{E/F}\Gm
 \longrightarrow T_{M,*}
 \xrightarrow{\ r\mapsto r/\iota(r)\ }K_T^1
 \longrightarrow1,
\end{equation}
where $K_T^1$ is the kernel of the norm from $K_T$ to $F$.
The map on $F$-points is surjective.  The quotient is split in
the first of the three cases above and anisotropic in the other
two.  In the field case $K_T$ is the third quadratic subfield
   of $M$.
\end{lemma}

\begin{proof}
For $t=r/\iota(r)$ and $r^*r=\lambda\in F^\times$,
\[
 \iota(t)=t^{-1},\qquad
 t^*=\frac{\lambda/r}{\lambda/\iota(r)}=t^{-1}.
\]
Thus $\iota *(t)=t$ and $t$ has norm one in $K_T$.
The kernel consists of the elements fixed by $\iota$, namely
$E^\times$.  Over $\barF$ this is the surjective quotient
of a three-dimensional split torus by its two-dimensional
central torus; equivalently, the ratios of the two primitive
coordinates give every norm-one pair.  This proves algebraic
exactness.  Hilbert~90 gives
$H^1(F,\Res_{E/F}\Gm)=1$, so exactness on rational points
gives surjectivity there as well.  In the exchanged-involution
case $\iota *$ is componentwise conjugation, whose fixed
algebra is $F\times F$.  In the componentwise-involution
case it has fixed algebra isomorphic to $E$.  In the field
case its fixed field is the third quadratic subfield.
\end{proof}

\subsubsection{The splitting kernels at admitted points}

We now fix points in the Hermitian domains just constructed.
The following maps measure splitting inside their geometric orbits.

\begin{proposition}[Rank-zero splitting]\label{m2:rat:rankzero-kernel}
Let $x=x_{\mu,D}$ be an admitted rank-zero $F$-point with closed
geometric orbit. The kernel in \eqref{m2:rat:orbit-kernel} is trivial.
For regular $D$ it is identified with $\ker\Delta$ for
$T_{E[D],*}\subset G_2$; for scalar $D$ it is the scalar
conjugacy problem.
\end{proposition}

\begin{proof}
The rational normalization in \eqref{m2:rat:rankzero-representative}
identifies rational orbits with $G_2(F)$-conjugacy classes of $D$.
Over $\barF$ the same normalization identifies the geometric
problem with conjugacy of its first loop component, with fixed
$\mu$. Thus two rational points in this geometric orbit reduce
to geometrically conjugate $D$-blocks. For regular $D$, apply
the cocycle construction of \cref{m2:rat:orbit-correspondence} to
the conjugation action of $G_2$; its stabilizer is $T_{E[D],*}$.
The splitting set is therefore
\begin{equation}\label{m2:rat:rankzero-regular-kernel}
 \ker[H^1(F,T_{E[D],*})\longrightarrow H^1(F,G_2)]
 =\ker\Delta=\{1\}
\end{equation}
by all three cases of \cref{m2:rat:torus-cohomology}.
Conversely, suppose
$hDh^{-1}=D'$ with $h\in G_2(F)$. Choose
$v=a_{\lambda(h)}$ and $u=a_\mu v a_\mu^{-1}$. Then
$\lambda(u)=\lambda(v)=\lambda(h)$ and
$(u,v,h)$ carries $x_{\mu,D}$ to $x_{\mu,D'}$.
This proves injectivity for the fiber-product stabilizer.

If $D=zI$, its conjugacy orbit is a point, so the same normalization
already proves uniqueness. Equivalently, the multiplier sequence
for $G_2\times_\lambda G_2$ has kernel $U_{1,1}^2$ and the
section $(a_s,a_s)$. The scalar-similarity argument proving
\eqref{m2:rat:H1K} gives $H^1(F,G_2\times_\lambda G_2)=Q_F^2$.
The three vertex forms of \eqref{m2:rat:rankzero-stabilizer} then give
\begin{equation}\label{m2:rat:rankzero-scalar-H1}
 (q_1,q_2)\longmapsto(q_1,q_1,q_2),
\end{equation}
an injective map. Conjugation of the first representation by
$a_\mu$ changes only the initial form, not its relative determinant.
\end{proof}

\begin{proposition}[Dense splitting]\label{m2:rat:dense-kernel}
Let $g=g_{\mu,D;B,C}$ be an admitted dense $F$-point with
closed geometric orbit. Its rational splitting kernel is trivial
in each of the absolutely irreducible, scalar, split-adjoint
toral, and anisotropic-adjoint toral cases.
\end{proposition}

\begin{proof}
The stabilizer embedding \eqref{m2:rat:dense-stabilizer}
allows us to treat the cases separately. For the absolutely
irreducible case its stabilizer is
$\Res_{E/F}\Gm$, whose $H^1$ is trivial by Hilbert~90.
If $D$ is scalar, the three two-dimensional representations are
conjugate copies of the defining $G_2$-representation. Pulling
back their forms gives scalar multiples of the hyperbolic form;
twisting changes the three relative determinant classes by the
same $q$. Thus
\begin{equation}\label{m2:rat:dense-scalar-H1}
 H^1(F,G_2)=Q_F\longrightarrow Q_F^3,
 \qquad q\longmapsto(q,q,q)
\end{equation}
has trivial kernel.

For regular toral $D$, the three representations are rank-one
$M$-modules, $M=E[D]$, with the same involution $*$: the
commuting Gram operators in \eqref{m2:rat:dense-Gram} do not change
the adjoint on $M$. Their trace-form coefficients may differ.
Twisting by $c\in L^\times$, however, changes each determinant
by the same ratio $\N_{L/F}(c)$, independently of that initial
coefficient, as proved in \cref{m2:rat:torus-cohomology}. Hence
\begin{equation}\label{m2:rat:dense-toral-H1}
 H^1(F,T_{M,*})\longrightarrow Q_F^3,
 \qquad[c]\longmapsto(\Delta[c],\Delta[c],\Delta[c]).
\end{equation}
For the exchanged-product involution the source is trivial.
For the componentwise-product involution and the field case,
the injectivity of $\Delta$ gives a trivial kernel. This covers
both anisotropic-adjoint forms as well as the split-adjoint form.
\end{proof}

\paragraph{Rank-one kernels.}

The factor tori of \cref{m2:rat:rankone-stabilizers} fit into the
algebraic norm-ratio sequence
\begin{equation}\label{m2:rat:Tm-sequence}
 1\longrightarrow T_m\longrightarrow(\Res_{E/F}\Gm)^m
 \xrightarrow{\psi}\Gm^{m-1}\longrightarrow1,
 \quad\psi(z)=\left(\frac{\N(z_1)}{\N(z_0)},\ldots,
                    \frac{\N(z_{m-1})}{\N(z_0)}\right).
\end{equation}
After splitting, independent primitive coordinates realize
every norm ratio, proving algebraic exactness. On $F$-points
the image is $\N(E^\times)^{m-1}$: each ratio is a norm, and
each tuple of norms is obtained with $z_0=1$. Hilbert~90
annihilates $H^1$ of the middle torus, so its rational cokernel is
\begin{equation}\label{m2:rat:H1Tm}
 H^1(F,T_m)=Q_F^{m-1}\simeq Q_F^m/\diag(Q_F).
\end{equation}
The last quotient describes changing the Hermitian coefficient
of each stable factor modulo norms and common scalar similarity.

The common determinant rule is as follows. A factor with
coefficient class $[c_i]$ and dimension $d_{vi}$ at vertex $v$
contributes $[c_i]^{d_{vi}}$ to that vertex's determinant ratio.
Orthogonality from \cref{m2:rat:selfadjoint-idempotents} makes the
whole ratio the product $\prod_i[c_i]^{d_{vi}}$.
Since every vertex has dimension two, simultaneous scaling of
all coefficients multiplies this product by a square and thus
does not change it in $Q_F$. The rule therefore descends to
the quotient in \eqref{m2:rat:H1Tm}.

\begin{proposition}[The two rank-one splitting maps]
\label{m2:rat:rankone-kernel}
At every admitted rank-one $F$-point with closed geometric
orbit, the stabilizer map to $H^1(F,K)=Q_F^3$ is
\begin{align}
 H^1(F,T_3)&\longrightarrow Q_F^3,&
 (\alpha,\beta)&\longmapsto(\alpha,\alpha,\beta),
       \label{m2:rat:T3-H1}\\
 H^1(F,T_2)&\longrightarrow Q_F^3,&
 \alpha&\longmapsto(\alpha,\alpha,1).
       \label{m2:rat:T2-H1}
\end{align}
Both maps have trivial kernel.
\end{proposition}

\begin{proof}
For $T_3$ use coefficient classes $(c_r,c_a,c_b)$ on
$(\mathrm R,\mathrm E,\mathrm Q)$ and set
$\alpha=[c_rc_a]$, $\beta=[c_rc_b]$. These are coordinates
on $Q_F^3/\diag(Q_F)$: because $Q_F$ has exponent two,
the product classes equal the ratio classes, and simultaneous
rescaling can make $[c_r]=1$. The vertex factor contents in
the $(u,v,h)$ order are $(r,a),(r,a),(r,b)$ by
\eqref{m2:rat:rankone-embeddings}. The determinant rule gives
\eqref{m2:rat:T3-H1}. A neutral image forces
$\alpha=\beta=1$, which is the neutral quotient class.

For $T_2$ the factors are $(\mathrm H,\mathrm E)$ with
classes $(c_r,c_a)$, and $\alpha=[c_rc_a]$ is the single
quotient coordinate. The contents are $(r,a),(r,a),(r,r)$.
The middle space belongs twice to the \emph{same} stable
factor, so its determinant ratio is $[c_r^2]=1$. This gives
\eqref{m2:rat:T2-H1}. A neutral image forces $\alpha=1$.
The argument uses only the factor dimensions and their
selfadjoint idempotents, so includes nilpotent and repeated-root
stable loops and every parameter in both Hermitian domains.
\end{proof}

\subsubsection{Closedness, uniqueness, and transpose}

The group $K$ is connected reductive: its kernel in
\eqref{m2:rat:ambient-sequence} is a product of connected
unitary groups, and the quotient is a torus.  Its
unipotent radical maps trivially to that torus and would
be a connected normal unipotent subgroup of the kernel,
so it is trivial.  The variety $G_4$ is affine.
The closed-orbit comparison theorem for reductive
actions over local fields
\cite[Theorem~2.3.8]{AGDescent} consequently identifies
analytically closed $K(F)$-orbits with geometrically
closed orbits at their rational points.

\begin{proof}[Proof of \cref{m2:rat:classification}]
Fix $x\in\mathscr O(F)$. The geometric classification
\cref{m2:geo:classification} exhausts the possibilities. The
rank-zero normalization, \cref{m2:rat:dense-param}, and
\cref{m2:rat:rankone-param} recover every admitted rational point;
their converses establish precisely the existence domains in
the statement. We apply \cref{m2:rat:orbit-correspondence} only
after this rational base point has been fixed. The stabilizers
are \eqref{m2:rat:rankzero-stabilizer},
\eqref{m2:rat:dense-stabilizer}, and
\eqref{m2:rat:rankone-embeddings}.
For the rank-zero, dense, and rank-one families, respectively,
\cref{m2:rat:rankzero-kernel,m2:rat:dense-kernel,m2:rat:rankone-kernel}
prove that the stabilizer kernel is trivial throughout these
domains. Thus \eqref{m2:rat:orbit-kernel} gives
\eqref{m2:rat:unique-orbit} in every row. The reductivity and
analytic closedness comparison just established gives the classification
of analytically closed $K(F)$-orbits.
\end{proof}

\begin{theorem}\label{m2:rat:transpose}
Every closed rational $K(F)$-orbit on $G_4(F)$ is
preserved by ordinary transpose.  For each of its
points there is $k_x\in K(F)$ with $k_x\cdot x^t=x$.
Its extended orbit under $\Kt(F)$ is the same
single $K(F)$-orbit.  If $\sigma_x=k_x\sigma$, then
\begin{equation}\label{m2:rat:extended-stabilizer}
 \Kt_x(F)=K_x(F)\sqcup K_x(F)\sigma_x,
 \qquad \sigma_x^2=k_x\alpha_K(k_x)\in K_x(F).
\end{equation}
\end{theorem}

\begin{proof}
\Cref{m2:geo:transpose} places $x$ and $x^t$ in the
same closed geometric orbit.  Both are rational,
and \eqref{m2:rat:unique-orbit} shows that this
geometric orbit has only one rational orbit.
Both also have the same multiplier by \cref{m2:red:transpose}.
This proves rational transpose stability. The explicit transporters in
Appendix~\ref{app:transporters} also satisfy the common-multiplier
condition and provide the choices used in the normal-space calculations.

The extended orbit is the union of $K(F)x$ and
$K(F)x^t$, so it is the same orbit.  Every odd
element fixing $x$ differs from $\sigma_x$ by
an element of $K_x(F)$, proving the coset
description.  Since $\sigma_x$ fixes $x$, its
square does as well, and its square belongs to the
even stabilizer, as expressed in \eqref{m2:rat:extended-stabilizer}.
Changing $k_x$ changes $\sigma_x$ by an even
stabilizer element.  This observation will allow
the normal calculations to use convenient
transporters without changing the invariant
distribution problem.
\end{proof}

\subsection{The normal representations}\label{m2:nor:section}

All normal spaces are over $F$. For a closed rational representative $x$,
write $S_x=K_x$ and choose $k_x=(u_x,v_x,h_x)\in K(F)$ with
$k_x\cdot x^t=x$ as in \cref{m2:rat:transpose}; put $\sigma_x=k_x\sigma$.
The odd action on $N_x$ is induced by this corrected transpose
$\sigma_x$, rather than by ordinary transpose alone. Another choice
of $k_x$ changes the action by an even stabilizer element.
We continue to write $\mathcal J=J\oplus J$.

\subsubsection{The tangent, quiver Ext, and corrected transpose}

For $\mu=\lam(x)$, left trivialize tangent vectors as $x\xi$.
Write $N_x$ for the normal space to $Kx$ in $G_4$, and $N_x^0$
for the reduced normal space in the fixed-multiplier fiber $X_\mu=\lam^{-1}(\mu)$.

\begin{proposition}[Normal spaces and scalar extension]\label{m2:nor:full}
For every rational closed-orbit representative $x$, the orbit differential is
\begin{equation}\label{m2:nor:beta}
 \beta_x(U,V,H)=\Ad(x^{-1})(U\oplus H)-(V\oplus H),
 \qquad d\lam(U)=d\lam(V)=d\lam(H).
\end{equation}
It has image in $\mathfrak u_4$, and
\begin{equation}\label{m2:nor:quotients}
 N_x^0=\frac{T_xX_\mu}{T_x(Kx)}
       =\frac{\mathfrak u_4}{\im\beta_x},\qquad
 N_x=\frac{\mathfrak{gu}_4}{\im\beta_x}
       =N_x^0\oplus F[I_4/2].
\end{equation}
Tangent spaces, $\Lie(K)$, and this differential commute with every
field extension $k/F$. If $k$ splits $E/F$, the comparison in
\cref{m2:geo:split-bridge} gives an equivariant identification
\begin{equation}\label{m2:nor:base-change}
 N_x^0\otimes_F k\simeq N^{\mathrm{lin}}_{x_1},
\end{equation}
where the right side is the normal space in the linear model's
determinant-open quiver variety. For $(u,v,h)\in S_x(F)$ and for
$\sigma_x$, respectively, the actions on \eqref{m2:nor:quotients} are induced by
\begin{align}
 \xi&\longmapsto\Ad(v\oplus h)\xi,\label{m2:nor:even}\\
 \Theta_x(\xi)&=
 \Ad\bigl((v_x\oplus h_x)(x^t)^{-1}\bigr)(\xi^t).
 \label{m2:nor:odd}
\end{align}
Both fix $I_4/2$. The odd square is the action of
$k_x\alpha_K(k_x)\in S_x(F)$; changing $k_x$ multiplies the
odd operator by an even stabilizer action.
\end{proposition}

\begin{proof}
Over $F[t]/(t^2)$ the equation for the Lie algebra is
$\xi^*+\xi=d\lam(\xi)I$, and
$\lam(x(1+t\xi))=\mu(1+t\,d\lam(\xi))$.
Thus $T_xX_\mu=x\mathfrak u_4$. Differentiating the orbit
map gives $(U\oplus H)x-x(V\oplus H)$, proving
\eqref{m2:nor:beta}. Its two terms have the same multiplier differential.
Its kernel is the differentiated stabilizer equation, hence $\Lie(S_x)$.
Algebraic groups and their stabilizers are smooth in characteristic zero;
the orbit map is separable, so its differential image is $T_x(Kx)$.

The tangent and Lie spaces are kernels of these $F$-linear matrix
equations. Tensoring gives their equations over $E\otimes_F k$, hence
\[
 T_xX_\mu\otimes_F k\simeq T_{x_k}(X_\mu)_k,\qquad
 \Lie(K)\otimes_F k\simeq\Lie(K_k).
\]
The differential tensors to $\beta_{x_k}$; flatness preserves its
kernel, image, and cokernel. Stabilizers commute with extension by
their equalizer equations.

When $k$ splits $E/F$, send $\xi$ to the additive tangent $x_1\xi_1$.
The equation $\xi^*+\xi=0$ determines its second component uniquely
from the unrestricted first component. Under this isomorphism
\eqref{m2:nor:beta} becomes
\begin{equation}\label{m2:nor:split-orbit-map}
 d\rho(L,Q_0,V)=
 (LA-AV,\ LB-BQ_0,\ Q_0C-CV,\ Q_0D-DQ_0),
\end{equation}
where $(L,Q_0,V)=(U_1,H_1,V_1)$ are unrestricted. The extra multiplier
Lie direction acts trivially. This is the quiver differential of
\cref{m1:lem:hom-ext,m1:eq:orbit-differential}. Fixing $\mu$ imposes
no condition on $\det x_1$, and the determinant open has the full
arrow-space tangent. Its cokernel proves \eqref{m2:nor:base-change}
equivariantly.

For an even element use $(u\oplus h)x=x(v\oplus h)$ to obtain
\eqref{m2:nor:even}. Transpose takes $x(1+t\xi)$ to
$(1+t\xi^t)x^t$. The transporter identity
$(u_x\oplus h_x)x^t=x(v_x\oplus h_x)$ then gives
\[
 (u_x\oplus h_x)(1+t\xi^t)x^t(v_x\oplus h_x)^{-1}
 =x(1+t\Theta_x(\xi)).
\]
This proves \eqref{m2:nor:odd}; the factor $(x^t)^{-1}$ converts
back to the left trivialization. It also verifies the split odd
compatibility explicitly: writing $y=x_1\xi_1$, and $a_1,b_1$
for the first components of $u_x\oplus h_x,v_x\oplus h_x$, one has
\begin{equation}\label{m2:nor:split-odd}
 x_1(\Theta_x\xi)_1=a_1y^tb_1^{-1}.
\end{equation}
Thus the split operator is the corrected point-fixing additive
transpose of the linear model and preserves the orbit tangent.
The semidirect product gives its square. If $k'_x$ is another
transporter, $s=k'_xk_x^{-1}\in S_x(F)$ gives
$\Theta'_x=\rho(s)\Theta_x$. After moving to a geometric representative,
the same comparison holds with its canonical linear transporter;
the extra multiplier torus acts trivially.

Finally $d\lam(I_4/2)=1$, while both \eqref{m2:nor:even} and
\eqref{m2:nor:odd} fix this central symmetric matrix. It splits the
multiplier differential equivariantly, proving the normal-space
splitting.
\end{proof}

The Hom--Ext complex of \cref{m1:lem:hom-ext}, applied to the
relation-free quiver with its loop, combines with
\eqref{m2:nor:base-change} to give
\begin{equation}\label{m2:nor:hom-ext}
 N_x^0\otimes_F k\simeq
 \operatorname{Ext}^1_{\mathscr Q}(V_{x_1},V_{x_1}).
\end{equation}
After splitting, \eqref{m1:eq:multiplicity-spaces} gives the
automorphism action on a decomposition into pairwise nonisomorphic
absolutely stable factors: each cross-Ext has target scalar divided
by source scalar. All arrows and the loop may deform.
As a dimension check,
\begin{equation}\label{m2:nor:dimension-check}
 \dim_F N_x^0=16-(13-\dim S_x)=3+\dim S_x.
\end{equation}

\subsubsection{Normal representations for dense orbits}

\begin{lemma}\label{m2:rat:dense-smooth}
On a fixed multiplier fiber, the map from the dense locus to
the $D$-coordinate is a smooth submersion onto its admissible
domain.  Its fibers are orbits of the subgroup of $K$ with
$h=1$.  Consequently its differential identifies the normal-space
quotient to a $K$-orbit with
\[
 \Mat_2(E)/\{[Z,D]:Z\in\Lie(G_2)\}.
\]
\end{lemma}

\begin{proof}
The fiber assertion follows from \cref{m2:rat:dense-param}. The
admissible domain is open on $F$-points: nondegeneracy is open and
determinant norm classes are locally constant on $F^\times$.
We verify the submersion and its tangent kernel explicitly.
For any nondegenerate Hermitian form $H$, the Gram derivative at $b$
on the increment $bZ$ is
\[
 Z^\dagger(b^\dagger Hb)+(b^\dagger Hb)Z.
\]
For a prescribed Hermitian increment $Y$, the choice
$Z=\tfrac12(b^\dagger Hb)^{-1}Y$ gives exactly $Y$.
Prescribe any $\dot D\in M_2(E)$ and differentiate
\eqref{m2:rat:dense-Gram}. The first equation in
\eqref{m2:rat:dense-factorizations} is solved by
$\dot B=\tfrac12 BQ_D^{-1}\dot Q_D$.
For the second, the right side of the $C$ Gram derivative is
$-C^\dagger\dot R_DC$, a Hermitian matrix; explicitly one may take
\[
 \dot C=\tfrac12 C(\mu J)^{-1}(-C^\dagger\dot R_DC).
\]
Define $\dot A$ by differentiating $A=BX_DC$. This solves all
linearized block equations, and the implicit function theorem over
$F$ gives local smooth Gram choices as $D$ varies.

If $\dot D=0$, then $\dot Q_D=\dot X_D=\dot R_D=0$.
The first linearized Gram equation is equivalent to
$\dot B=UB$ with a unique $U\in\mathfrak u_2$; the second is
equivalent to $\dot C=-CV$ with a unique $V\in\mathfrak u_2$.
Thus the tangent kernel is
\begin{equation}\label{m2:rat:dense-tangent-kernel}
 (\dot A,\dot B,\dot C,\dot D)=(UA-AV,UB,-CV,0),
 \qquad U,V\in\mathfrak u_2,
\end{equation}
the tangent of the $(u,v,1)$ gauge orbit. After quotienting by it,
the residual differential is $\dot D=[H,D]$ with
$H\in\mathfrak{gu}_2$. Every such $H$ is allowed in $\Lie(K)$,
by choosing $U,V$ with its multiplier differential. This proves
the asserted normal quotient.
\end{proof}

\begin{lemma}\label{m2:nor:dense-quotient}
On the entire admitted dense domain of \cref{m2:rat:dense-param},
\begin{equation}\label{m2:nor:dense-normal}
 N_x^0\simeq M_2(E)/[\mathfrak{gu}_2,D].
\end{equation}
The even action is $[Y]\mapsto[hYh^{-1}]$ for
$h\in Z_{G_2}(D)$, and the corrected odd map is
\begin{equation}\label{m2:nor:dense-odd}
 [Y]\longmapsto[h_xY^th_x^{-1}].
\end{equation}
\end{lemma}
\begin{proof}
\Cref{m2:rat:dense-smooth} proves the unitary Gram submersion and
identifies its tangent kernel with the $(u,v,1)$ gauge tangent.
The remaining differential is $[H,D]$, giving the quotient.
The embedding \eqref{m2:rat:dense-stabilizer} gives the even action.
Transpose and its transporter send the lower-right additive block
to $h_xD^th_x^{-1}$, whose derivative is \eqref{m2:nor:dense-odd}.
The norm-class conditions define open subsets, so the tangent
calculation applies throughout the admitted domain.
\end{proof}

With $\delta,P,R$ and $\Psi$ as in \cref{m2:red:unitary-pair}, we write
\begin{equation}\label{m2:nor:pair-coordinate}
 P^{-1}DP=A_0+\delta B_0,\qquad A_0,B_0\in M_2(F).
\end{equation}
The factorization $h=zPgP^{-1}$ gives the equivalence
\[
 hD=Dh\quad\Longleftrightarrow\quad
 gA_0=A_0g\ \hbox{and}\ gB_0=B_0g.
\]
Every such $g\in\GL_2(F)$ lifts, so these matrices give the
effective rational stabilizer action.
The Lie equation in the $P$ basis is
$\bar Z^tJ_0+J_0Z=cJ_0$. Writing $Z=X+\delta Y$ with real matrices
gives $\tr X=c$ and $Y=(\tr Y)I/2$; hence
$P^{-1}\mathfrak{gu}_2P=\mathfrak{gl}_2(F)+EI$.
Its central part has zero commutator. Therefore $\Psi$ induces
\begin{equation}\label{m2:nor:pair-normal}
 N_x^0\simeq
 \frac{M_2(F)^2}
 {\{([Z,A_0],[Z,B_0]):Z\in\mathfrak{gl}_2(F)\}}.
\end{equation}
Using $\mathfrak{sl}_2(F)$ gives the same image. The algebra identities
\cref{m2:rat:real-pair} and \eqref{m2:rat:dense-algebra} show that closedness
of $x$ makes this real pair geometrically semisimple.

Choose the symmetric intertwiner $H$ in
\cref{m2:rat:dense-transporter}, so $A_0H=HA_0^t$, $B_0H=HB_0^t$,
and $h_x=P(HR^{-1})P^{-1}$. The uncorrected transpose in
\eqref{m2:simconj:transpose} contributes $R$; conjugation by $h_x$
cancels it. Thus the corrected pair operator is
\begin{equation}\label{m2:nor:pair-odd}
 (Y_1,Y_2)\longmapsto(HY_1^tH^{-1},HY_2^tH^{-1}).
\end{equation}
Other rational transporters add the stabilizer action just identified.

Choose scalar translations $(A_0+sI,B_0+tI)$ in the linear
determinant open as in \cref{m2:geo:pair-transpose}. They preserve the
generated algebra, centralizer, commutator image, and intertwiner
equations. Their differential is the identity and commutes with
\eqref{m2:nor:pair-odd}. Thus \cref{m1:eq:dense-normal} identifies the
quotient, even action, and corrected odd operator with the normal
representation for a dense orbit in the linear model, using $H=H_0^{-1}$ in \cref{m1:lem:dense-transpose}.

\paragraph{Absolute and scalar pairs}

For an absolutely irreducible pair, \cref{m1:prop:abs-irred-normal}
identifies the quotient by its five invariant differentials and
proves pointwise odd fixation. Hence
\begin{equation}\label{m2:nor:dense-simple}
 S_x=\Res_{E/F}\Gm,\qquad N_x^0=F^5_{\triv},\qquad\Theta_x=1.
\end{equation}

Equivalently, use the differentials of
\begin{equation}\label{m2:nor:five-invariants}
 \tr A_0,\quad\tr B_0,\quad\tr(A_0^2),\quad
 \tr(A_0B_0),\quad\tr(B_0^2).
\end{equation}
The identity $\tr(C^2)=(\tr C)^2-2\det C$ gives an invertible change
from the trace/determinant differentials. For a scalar pair,
\cref{m1:subsec:dense-normals} gives $M_2(F)^2$ with simultaneous
conjugation and componentwise transpose, hence $N_x^0=M_2(E)$ with
$G_2$-conjugation. For $D=d_cI$, \cref{m2:rat:dense-transporter} permits
$h_x=1$, $u_x=BC^{-t}$, $v_x=C^{-1}B^t$, all of multiplier one.
Then \eqref{m2:nor:dense-odd} is $Y\mapsto Y^t$; its factor $R$ in
pair coordinates is an even conjugation.

\paragraph{Toral pairs and the root field}

For a toral pair let $L_0=F[A_0,B_0]$, a quadratic etale algebra.
Its centralizer is $L_0^\times$. In the transported algebra
$M=E\otimes_F L_0$, the $E$-linear involution $\iota$ restricts
to the nontrivial automorphism of $L_0$. The adjoint formula in
\cref{m2:rat:real-pair} restricts to
$* =\overline{\phantom z}_E\otimes\iota$:
on a quadratic matrix algebra, adjugation is its nontrivial
automorphism. Consequently
\begin{equation}\label{m2:nor:root-field}
 K_T=M^{\iota *}=L_0,
\end{equation}
as embedded $F$-algebras. This identifies the real-pair root field
with \cref{m2:rat:adjoint-torus}.

If $L_0=F\times F$, the normal-space and transpose formulas
for the split case in \cref{m1:eq:dense-split-normal} give
\begin{equation}\label{m2:nor:dense-split}
 N_x^0=F^4_{\triv}\oplus F_\alpha\oplus F_{\alpha^{-1}},
 \qquad\Theta_x(z,u,v)=(z,v,u).
\end{equation}
The four self-Ext directions are odd-fixed by that result. In a
real simultaneous eigenbasis write $g=\diag(t_1,t_2)$ in the
factorization of $h$. The character is
$\alpha(h)=t_1/t_2$; it is independent of the common-scalar overlap.
The lift of $\diag(q,1)$ realizes every $q\in F^\times$.

For a field $L_0$, use the quotient and action in
\cref{m1:eq:dense-nonsplit-quotient,m1:eq:dense-nonsplit-action}:
if $A_0=m_a$, $B_0=m_b$, put $r_a=\bar a-a$, $r_b=\bar b-b$.
At least one is nonzero, and the root quotient is
\begin{equation}\label{m2:nor:anis-root-coordinate}
 L_0^2/L_0(r_a,r_b)\xrightarrow{\ \ell\ }L_0,
 \qquad\ell(y_1,y_2)=r_by_1-r_ay_2.
\end{equation}
The cited calculation gives multiplication $w\mapsto(t/\bar t)w$
and corrected transpose $w\mapsto-\bar w$, fixing the other four
coordinates. Bars here mean the $L_0/F$ involution. For nonzero
trace-zero $\omega\in L_0$, the coordinate $w=\omega w'$ commutes
with root multiplication and satisfies
$-\overline{\omega w'}=\omega\bar w'$. Hence
\begin{equation}\label{m2:nor:dense-anis}
 N_x^0=F^4_{\triv}\oplus K_T,\qquad
 \Theta_x(z,w')=(z,\bar w').
\end{equation}
Every $t\in L_0^\times$ lifts through $PgP^{-1}$, and Hilbert~90
makes $t/\bar t$ range over $K_T^1$, as in
\cref{m2:rat:adjoint-torus}. This includes $K_T\simeq E$ and the
biquadratic case, where $K_T$ is the third quadratic field.

\subsubsection{Rank-one normal representations and rational descent}

We use \eqref{m2:nor:hom-ext} over a splitting field $k$. For the
decomposable factors $(\mathsf E,\mathsf R,\mathsf Q)$, of dimensions
$(1,0,1),(1,1,1),(0,1,0)$, denote their scalars by $(a,r,b)$.
The vertex actions are $p=v=\diag(r,a)$, $q=\diag(r,b)$.
The linear Ext table \eqref{m1:eq:decomposable-ext-table} gives
\begin{equation}\label{m2:nor:dec-geometric}
 (N_x^0)_k=k^3_{\triv}\oplus k_\alpha\oplus k_{\alpha^{-1}}
                         \oplus k_\beta\oplus k_{\beta^{-1}},
 \qquad \alpha=r/b,\quad\beta=a/b.
\end{equation}
The positive lines are $\Ext^1(\mathsf Q,\mathsf R)$ and
$\Ext^1(\mathsf Q,\mathsf E)$; the trivial part consists of two
$\mathsf R$ self-Ext classes and one $\mathsf Q$ self-Ext class.
By \cref{m1:eq:decomposable-fixed-coordinates,m1:eq:decomposable-transpose},
the canonical corrected operator $\tau_0$ fixes these three classes,
swaps each opposite pair, inverts the effective torus, and squares
to one. This holds on the entire domain $\omega\ne0$, $\zeta\nu\ne1$,
including admitted $\zeta=0$ or $\nu=\omega$.

For the stable system $\mathsf E\oplus\mathsf H$ with scalars
$(a_s,r)$, the vertex actions are $p=v=\diag(r,a_s)$, $q=rI$.
The table \eqref{m1:eq:stable-ext-table} gives
\begin{equation}\label{m2:nor:stable-ext-table}
 (N_x^0)_k=k^4_{\triv}\oplus k_\eta\oplus k_{\eta^{-1}},
 \qquad\eta=a_s/r.
\end{equation}
Here the $\eta$ line is $\Ext^1(\mathsf H,\mathsf E)$ and the
trivial part is $\Ext^1(\mathsf H,\mathsf H)$. The canonical
transporter $(I,H_0^{-1},I)$ in vertex order $(p,q,v)$ induces
\begin{equation}\label{m2:nor:stable-odd-geometric}
 (\dot A,\dot B,\dot C,\dot D)\longmapsto
 (\dot A^t,\dot C^tH_0,H_0^{-1}\dot B^t,H_0^{-1}\dot D^tH_0).
\end{equation}
The parameter derivatives in $a,b_0,d_0,s_0$ form a basis of the
self-Ext part and are fixed modulo orbit tangents by
\cref{m1:lem:stable-parameter-basis,m1:eq:stable-parameter-derivatives,m1:eq:stable-ambient-transpose}:
differentiating the parameter-dependent transporter contributes an
orbit tangent. The same operator squares to one, inverts $\eta$,
and exchanges its root lines for every $\det H_0\ne0$, including
nilpotent and repeated-root loops. Compatibility with the rational normal representation
follows from \eqref{m2:nor:split-orbit-map}--\eqref{m2:nor:split-odd}.

\begin{lemma}[Descent of a multiplicity-one opposite pair]
\label{m2:nor:root-pair-descent}
Let $T$ be an $F$-torus split over a quadratic field $L/F$. Suppose an
$F$-representation $W$ has
$W_L=L_\chi\oplus L_{\chi^{-1}}$, two distinct one-dimensional
weight spaces of multiplicity one, exchanged by Galois. Let an
$F$-linear odd operator exchange these lines and suppose its square
is the identity on this effective module. Then $W$ is the underlying
$F$-plane of $L$, with its corresponding norm-one multiplication
character. The odd map is $z\mapsto c\bar z$, with $c\bar c=1$,
and an $L$-scalar coordinate change makes it ordinary conjugation.
\end{lemma}
\begin{proof}
Choose $e_\chi$ and its Galois conjugate $\bar e_\chi$ as the two
bases. The rational vectors are
$ze_\chi+\bar z\bar e_\chi$. Galois sends $\chi$ to $\chi^{-1}$,
so $\chi(s)\overline{\chi(s)}=1$ for $s\in T(F)$ and the action
in this coordinate is multiplication by $\chi(s)$. Exchanging the
lines and commuting with descent makes the odd map $c\bar z$.
Its assumed square one gives $c\bar c=1$ before Hilbert~90 is
used. Write $c=b/\bar b$ and $z=bw$; the map becomes $w\mapsto\bar w$.
This change commutes with the multiplication action.
\end{proof}
The algebraic stabilizers, with embeddings
\eqref{m2:rat:rankone-embeddings}, are
\begin{equation}\label{m2:nor:rational-rankone-tori}
 \begin{split}
 T_3&=\{(r,a,b)\in(\Res_{E/F}\Gm)^3:N(r)=N(a)=N(b)\},\\
 T_2&=\{(r,a)\in(\Res_{E/F}\Gm)^2:N(r)=N(a)\}.
 \end{split}
\end{equation}
They split over $E$. Define the norm-one characters
\begin{equation}\label{m2:nor:rational-characters}
 \chi_\alpha(r,a,b)=r/b,\qquad
 \chi_\beta(r,a,b)=a/b,\qquad\eta(r,a)=a/r.
\end{equation}
For such a character, $E_\chi$ denotes the underlying $F$-plane
of $E$ with multiplication by $\chi$.

\begin{proposition}\label{m2:nor:rankone-rational}
Throughout the admitted rational domains,
\begin{equation}\label{m2:nor:rankone-forms}
 \begin{array}{c|c|c}
 \text{family}&N_x^0&\Theta_x\\\hline
 \text{decomposable}&F^3_{\triv}\oplus E_{\chi_\alpha}
                 \oplus E_{\chi_\beta}&(z,u,v)\mapsto(z,\bar u,\bar v)\\
 \text{stable system}&F^4_{\triv}\oplus E_\eta
                 &(z,u)\mapsto(z,\bar u)
 \end{array}
\end{equation}
in suitable $F$-linear coordinates. The effective maps
$T_3(F)\to(E^1)^2$ and $T_2(F)\to E^1$ given by
\eqref{m2:nor:rational-characters} are surjective.
\end{proposition}
\begin{proof}
The common-norm equations give
$\overline{r/b}=b/r$, $\overline{a/b}=b/a$, $\overline{a/r}=r/a$.
Thus Galois exchanges each opposite pair of the characters in
\eqref{m2:nor:dec-geometric} and \eqref{m2:nor:stable-ext-table}.
The decomposable pairs are distinct and cannot be mixed by descent
or the odd map.

It remains to verify the odd hypotheses over $F$. By
\eqref{m2:nor:split-odd}, the rational transporter gives, after passage
to the geometric representative, $\rho(s)\tau_0$ for a stabilizer
element $s$. The preceding $\tau_0$ squares to one and inverts the
effective torus. Hence on the entire effective normal module
\[
 (\rho(s)\tau_0)^2=\rho(s)\rho(s^{-1})=1.
\]
The common scalar and the additional split multiplier act trivially.
Thus the rational operator is involutive and exchanges each pair.
With $L=E$, \cref{m2:nor:root-pair-descent} identifies each root plane
with $E_\chi$ and shows that its odd coefficient has norm one;
Hilbert~90 normalizes it to conjugation. The two decomposable
coordinate changes can be made independently.

The three or four even-trivial coordinates are odd-fixed by the
linear operator calculations above; stabilizer corrections act
trivially there, so this descends over $F$. Finally
$(r,a,b)=(q_1,q_2,1)$ realizes any $(q_1,q_2)\in(E^1)^2$ jointly,
and $(r,a)=(1,q)$ realizes any $q\in E^1$, all with common norm one.
\end{proof}

\subsubsection{The three rank-zero families}

We take the representative $x=a_\mu\oplus D$ from
\cref{m2:rat:classification}, where $a_\mu=\diag(1,\mu)$,
$\lam(D)=\mu$, and $D$ is semisimple.  Its stabilizer is
\[
 S_x=\{(a_\mu v a_\mu^{-1},v,h):
       (v,h)\in G_2\times_\lam Z_{G_2}(D)\}.
\]
For $Q\in M_2(E)$ put $Q^*=JQ^\dagger J$.  Every vector in
$\mathfrak u_4$ has the unique form
\begin{equation}\label{m2:nor:unitary-blocks}
 \xi=\begin{pmatrix}P_0&Q\\-Q^*&S_0\end{pmatrix},
 \qquad P_0,S_0\in\mathfrak u_2,
 \quad Q\in M_2(E).
\end{equation}

\begin{proposition}\label{m2:nor:zero-normal}
For the rank-zero representative above,
\[
 N_x^0\simeq M_2(E)\oplus\mathfrak c_D,
 \qquad
 \mathfrak c_D=
 \begin{cases}
 \mathfrak u_2,&D\text{ scalar},\\
 \mathfrak u_2\cap E[D],&D\text{ regular semisimple}.
 \end{cases}
\]
With $S_x$ identified as above, the stabilizer acts by
\begin{equation}\label{m2:nor:zero-action}
 (v,h)\cdot(Q,S_0)=(vQh^{-1},hS_0h^{-1}).
\end{equation}
Put
\[
 c_U=v_xa_\mu^{-t},\qquad c_W=h_xD^{-t}.
\]
Then
\begin{equation}\label{m2:nor:zero-odd}
 \Theta_x(Q,S_0)=
 \bigl(-c_UJ\bar QJc_W^{-1},\ c_WS_0^tc_W^{-1}\bigr),
 \qquad
 \lam(c_U)=\lam(c_W)=\frac{\lam(k_x)}\mu,
 \qquad c_WD^tc_W^{-1}=D.
\end{equation}
For regular $D$ the two-dimensional space $\mathfrak c_D$ is
pointwise fixed by both actions.  For scalar $D=d_cI$, the odd
operator may be taken, up to an even action, to be
\[
 (Q,S_0)\longmapsto(-J\bar QJ,S_0^t).
\]
\end{proposition}

\begin{proof}
At $a_\mu\oplus D$, the image in \eqref{m2:nor:beta} is block
diagonal, with blocks
\[
 a_\mu^{-1}Ua_\mu-V,
 \qquad D^{-1}HD-H.
\]
The first block fills $\mathfrak u_2$, by taking $V=H=0$ and
$U=a_\mu P_0a_\mu^{-1}$.  Conversely, after subtracting a common
central infinitesimal scalar from $U,V,H$, we may suppose that all
three have multiplier differential zero; that subtraction does
not change either of the two blocks just computed.  Consequently the image is exactly
\[
 \mathfrak u_2\oplus
 (\Ad(D^{-1})-1)\mathfrak u_2
\]
in the two diagonal blocks, and it has no off-diagonal part.

Semisimplicity of $D$ gives, after a splitting extension and then
descent to $F$,
\[
 \mathfrak u_2=
 \ker(\Ad(D)-1)\oplus\im(\Ad(D)-1).
\]
The kernel is $\mathfrak u_2$ for scalar $D$, and
$\mathfrak u_2\cap E[D]$ for regular $D$, since the matrix
centralizer is $E[D]$. In the latter case it is the Lie algebra of
the two-dimensional unitary centralizer. This proves the quotient;
\eqref{m2:nor:even} gives its even action.

Formula \eqref{m2:nor:odd} now has block-diagonal conjugating matrix
$c_U\oplus c_W$.  The top-right block of $\xi^t$ is
$-(Q^*)^t=-J\bar QJ$, giving \eqref{m2:nor:zero-odd}.
The multiplier equality follows from
$\lam(v_x)=\lam(h_x)$ and
$\lam(a_\mu)=\lam(D)=\mu$.  The lower-right transporter
equation is $h_xD^th_x^{-1}=D$; multiplying $h_x$ by $D^{-t}$
does not change this conjugation, so
$c_WD^tc_W^{-1}=D$ as claimed.  If $S_0\in E[D]$, write it
as $f(D)$ with coefficients in $E$.  Transpose leaves those scalar
coefficients unchanged, and therefore
$c_WS_0^tc_W^{-1}=f(D)=S_0$.  The centralizer also fixes this
space by conjugation.

If $D=d_cI$, the representative is symmetric and we can take
$k_x=1$.  Thus $c_U=a_\mu^{-1}$ and $c_W=d_c^{-1}I$.
The resulting leg map is $(d_ca_\mu^{-1})(-J\bar QJ)$.
Since $N_{E/F}(d_c)=\mu$, the matrix $d_ca_\mu^{-1}$ has
multiplier one.  Its left action, with right factor $1$, is an even
stabilizer action.  Removing it leaves the stated odd operator and
does not alter $S_0^t$.
\end{proof}

The scalar row has dimensions $(\dim S_x,\dim N_x^0)=(9,12)$,
and each regular row has $(7,10)$. In the regular rows the two
central directions are fixed, while the adjoint-torus type affects
the leg representation with the conditions in \eqref{m2:nor:zero-odd}.

\subsubsection{Classification of the rational normal representations}

\begin{theorem}\label{m2:nor:classification}
The nine rational closed-orbit families of \cref{m2:rat:classification}
have the algebraic $F$-stabilizers and reduced normal spaces
in the table below, with their indicated parameter-domain references.
Here $T=T_{M,*}$ is the torus of \cref{m2:rat:torus-cohomology}; its
split or anisotropic adjoint type and $K_T$ are those of
\cref{m2:rat:adjoint-torus}. The stabilizer embeddings in the table are \eqref{m2:rat:rankzero-stabilizer}, \eqref{m2:rat:dense-stabilizer},
and \eqref{m2:rat:rankone-embeddings}. Put
\[
 \kappa(Q)=-J\bar QJ,\qquad
 \kappa_x(Q)=c_U\kappa(Q)c_W^{-1},\qquad
 c_U=v_xa_\mu^{-t},\quad c_W=h_xD^{-t}.
\]
In both regular rank-zero rows the transporter conditions are
\begin{equation}\label{m2:nor:table-transport-conditions}
 \lam(c_U)=\lam(c_W)=\lam(k_x)/\mu,\qquad
 c_WD^tc_W^{-1}=D.
\end{equation}
Every $F^j$-summand in the table is fixed by both actions. In the action column,
only the other coordinates are written; the first line is even
and the second is odd. The last column gives
$(\dim S_x,\dim_F N_x^0)$.

{\small
\begin{longtable}{@{}L{.20\textwidth}L{.31\textwidth}L{.32\textwidth}c@{}}
\toprule
Family and domain & $S_x$; reduced normal $N_x^0$
 & Even action; corrected odd action & Dims.\\
\midrule
\endfirsthead
\toprule
Family and domain & $S_x$; reduced normal $N_x^0$
 & Even action; corrected odd action & Dims.\\
\midrule
\endhead
\bottomrule
\endfoot
Zero scalar,
$D=d_cI$, $N(d_c)=\mu$;
\eqref{m2:rat:rankzero-representative}
 & $G_2\times_\lam G_2$;
   $M_2(E)\oplus\mathfrak u_2$
 & $(Q,Z)\mapsto(vQh^{-1},hZh^{-1})$;
   $(Q,Z)\mapsto(a_\mu^{-1}\kappa(Q)d_c,Z^t)$
 & $9,12$\\[5pt]
Zero regular split-adjoint;
\eqref{m2:rat:rankzero-representative}
 & $G_2\times_\lam T$;
   $M_2(E)\oplus F^2_{\triv}$
 & $Q\mapsto vQh^{-1}$;
   $Q\mapsto\kappa_x(Q)$
 & $7,10$\\[5pt]
Zero regular anisotropic-adjoint;
\eqref{m2:rat:rankzero-representative}
 & $G_2\times_\lam T$;
   $M_2(E)\oplus F^2_{\triv}$
 & $Q\mapsto vQh^{-1}$;
   $Q\mapsto\kappa_x(Q)$
 & $7,10$\\[5pt]
Dense absolute;
\eqref{m2:rat:dense-domain}
 & $\Res_{E/F}\Gm$;
   $F^5_{\triv}$
 & identity; identity
 & $2,5$\\[5pt]
Dense scalar;
\eqref{m2:rat:dense-domain}
 & $G_2$;
   $M_2(E)$
 & $Y\mapsto hYh^{-1}$;
   $Y\mapsto Y^t$ ($h_x=1$)
 & $5,8$\\[5pt]
Dense split-adjoint toral;
\eqref{m2:rat:dense-domain}
 & $T$;
   $F^4_{\triv}\oplus F_\alpha\oplus F_{\alpha^{-1}}$
 & $(u,v)\mapsto(\alpha(h)u,\alpha(h)^{-1}v)$;
   $(u,v)\mapsto(v,u)$
 & $3,6$\\[5pt]
Dense anisotropic-adjoint toral;
\eqref{m2:rat:dense-domain}
 & $T$;
   $F^4_{\triv}\oplus K_T$
 & $w\mapsto\chi_T(h)w$;
   $w\mapsto\bar w$ over $K_T/F$
 & $3,6$\\[5pt]
One decomposable;
\eqref{m2:rat:decomposable-domain},
\eqref{m2:rat:decomposable-norms}
 & $T_3$;
   $F^3_{\triv}\oplus E_{\chi_\alpha}\oplus E_{\chi_\beta}$
 & $(u,v)\mapsto(\chi_\alpha u,\chi_\beta v)$;
   $(u,v)\mapsto(\bar u,\bar v)$
 & $4,7$\\[5pt]
One stable system;
\eqref{m2:rat:stable-domain},
\cref{m2:rat:rankone-param}
 & $T_2$;
   $F^4_{\triv}\oplus E_\eta$
 & $u\mapsto\eta u$;
   $u\mapsto\bar u$
 & $3,6$
\end{longtable}}

Here $\alpha(T(F))=F^\times$, and
$\chi_T:T(F)\twoheadrightarrow K_T^1$, $h\mapsto h/\iota(h)$,
is the root ratio in \cref{m2:rat:adjoint-torus}. The rank-one maps
\eqref{m2:nor:rational-characters} are surjective, jointly for two
roots. Both allowed anisotropic involutive-algebra forms are included.
The scalar zero-row operator uses $k_x=1$; removing the even action
$(d_ca_\mu^{-1},1)$ gives $(\kappa(Q),Z^t)$ and retains the entire
moving $\mathfrak u_2$. Other rows use the coordinates and rescalings
proved above.

For every row the normal space and extended stabilizer are
\begin{equation}\label{m2:nor:full-classification}
 N_x=N_x^0\oplus F[I_4/2],\qquad
 \widetilde S_x(F)=S_x(F)\sqcup S_x(F)k_x\sigma.
\end{equation}
The multiplier line is pointwise fixed. The sign $\eps$ is $1$
on $S_x(F)$ and $-1$ on its odd coset. A new transporter adds an
even stabilizer action; a new point of the same rational orbit
transports the entire representation by its rational group element.
The absence of sign-equivariant distributions on $N_x^0$
therefore implies the same on $N_x$.
\end{theorem}
\begin{proof}
\Cref{m2:nor:zero-normal} gives the three unitary rank-zero rows with
their transporter conditions. \Cref{m2:nor:dense-quotient} and
\eqref{m2:nor:dense-simple}--\eqref{m2:nor:dense-anis} give the four dense
rows through the rational matrix-pair comparison. \Cref{m2:nor:rankone-rational}
gives the remaining two. Their stabilizers and admitted domains are
the unchanged ones in \cref{m2:rat:classification}, which proves
exhaustiveness. All nine dimensions agree with
\eqref{m2:nor:dimension-check}. Finally \cref{m2:nor:full} proves the
normal-space splitting and transporter comparison. Projection to its
fixed multiplier line has complete fibers isomorphic to $N_x^0$;
\cref{m2:tools:localization} proves the last assertion.
\end{proof}
